\documentclass[aspectratio=169,xcolor=dvipsnames]{beamer}
\usetheme{SimpleDarkBlue}
\setbeamertemplate{caption}[numbered] % Enables numbered captions

\useoutertheme[subsection=false]{smoothbars}
\useinnertheme{circles} % clean bullet style
\setbeamertemplate{enumerate items}[default]

% --- Colors for smoothbars ---
\setbeamercolor{section in head/foot}{bg=DarkBlue, fg=white}         % active section bar
\setbeamercolor{section in head/foot shaded}{bg=LightBlue, fg=white}   % inactive bars
\setbeamercolor{frametitle}{bg=DarkBlue, fg=white}                     % frame title bar

\usepackage{url} %this package should fix any errors with URLs in refs.
\usepackage[authoryear]{natbib}
\bibpunct[]{(}{)}{;}{a}{,}{,}

\usepackage{caption}
\usepackage{hyperref}
\usepackage{graphicx} % Allows including images
\usepackage{booktabs} % Allows the use of \toprule, \midrule and \bottomrule in tables

\usepackage{geometry}

% --- Macro to start a new appendix-like section ---
\newcommand{\startappendix}{
  \setcounter{framenumber}{0}  % reset numbering
  \setbeamertemplate{footline}{}  % clear first
  % Define footline to show ONLY current frame number
  \setbeamertemplate{footline}{
    \hfill\insertframenumber\hspace{2ex}\vspace{2ex}
  }
}

% ----------------------------------------------------------------------------------
%	TITLE PAGE
%----------------------------------------------------------------------------------

% The title

\title[short title]{The relation of ENSO with global mean temperature}

\author[]{Michael K. Tippett\inst{1} and Emily J. Becker\inst{2}}
\institute[]{\inst{1} Department of Applied Physics and Applied
  Mathematics, Columbia University, New York, New York \and %
  \inst{2} University of Miami Rosenstiel School for Marine, Earth,
  and Atmospheric Science, Miami, Florida}

\date{\today} % Date, can be changed to a custom date

% Define section numbering style
%\renewcommand{\thesection}{\arabic{section}}          

\setcounter{secnumdepth}{1} % Ensure section numbers are assigned

\AtBeginSection[]{
    \setbeamercolor{background canvas}{bg=blue!50} % Light blue background
    \begin{frame}
        \centering
        \Huge\textbf{\textcolor{white}{\thesection.\ \insertsectionhead}} % Display section number and name
      \end{frame}
    \setbeamercolor{background canvas}{bg=white} % Reset to default background color
}

% Automatically create a slide with the subsection name and number at the
% beginning of each subsection
\AtBeginSubsection[]{
    \setbeamercolor{background canvas}{bg=blue!50} % Light blue background
    \begin{frame}
        \centering
        \LARGE\textbf{\textcolor{white}{\thesection.\thesubsection\ \insertsubsectionhead}} % Display subsection number and name
    \end{frame}
    \setbeamercolor{background canvas}{bg=white} % Reset to default background color
}

% ----------------------------------------------------------------------------------
%	PRESENTATION SLIDES
%----------------------------------------------------------------------------------

\begin{document}
%TC:ignore

\begin{frame}
    % Print the title page as the first slide
    \titlepage
\end{frame}

% Introduction--why the study was undertaken?

% What question did you ask in your experiment? Why is it interesting?
% The introduction summarizes the relevant literature so that the
% reader will understand why you were interested in the question you
% asked. One to four paragraphs should be enough. End with a sentence
% explaining the specific question you asked in this experiment.

% First of all, explanation of the topic in the light of the current
% literature should be made in clear, and precise terms as if the
% reader is completely ignorant of the subject.

% 
% Then main topic of our manuscript, and the encountered problem should
% be analyzed in the light of the current literature.

% Finally, the last paragraphs of the ‘Introduction’ section should
% include the solution in which we will describe the information we
% generated, and related data.

%The Introduction section clarifies the motivation for the work
%presented and prepares readers for the structure of the paper.

% First, provide some context to orient those readers who are less
% familiar with your topic and to establish the importance of your
% work.  Second, state the need for your work, as an opposition
% between what the scientific community currently has and what it
% wants.  Third, indicate what you have done in an effort to address
% the need (this is the task).  Finally, preview the remainder of the
% paper to mentally prepare readers for its structure, in the object
% of the document.

% Introduction: What did you/others do? Why did you do it?

% What is the problem to be solved?
% Are there any existing solutions?
% Which is the best?
% What is its main limitation?
% What do you hope to achieve?
% 

%TC:endignore

% \begin{frame}{Key points}
%   \begin{enumerate}
%   \item The oft-touted three-month lag correlation of ENSO with global
%     mean temperature is wrong in two important ways.
%     \begin{itemize}
%     \item The relation is not causal.
%     \item Only true one month out of the year
%     \end{itemize}
%   \item Despite being unrecognized for decades, these truths are
%     simple consequences of well-known stuff: serial correlation and
%     seasonality of the ENSO atmospheric bridge.
%   \end{enumerate}
% \end{frame}

\section{Abstract}

\begin{frame}{Abstract}
  \begin{enumerate}\scriptsize
  \item ENSO is a predictable contribution to global temperature.
  \item Previous studies found that the lag correlation between ENSO and
    global mean temperature peaks at 3--4 months with ENSO leading.
  \item Here we examined this relation with updated data and attention
    to seasonality.
  \item We found weaker relations than previously reported and that
    accounting for seasonality limits the peak correlations between
    ENSO and global mean temperature to January and February.
  \item The seasonality of the relation between ENSO and global mean
    temperature is explained by the seasonality of ENSO itself and by
    the seasonality of the atmospheric bridge mechanisms through which
    ENSO-related circulation drives remote SST changes.
  \end{enumerate}
\end{frame}

\section{Introduction}
% \begin{frame}{Introduction outline}
%   \begin{enumerate}
%   \item The problem: What's the deal with ENSO and global
%     temperatures?
%   \item Previous work: ENSO and global mean temperature
%   \item Previous work: Atmospheric bridge
%   \item The gaps and how we filled them
%   \end{enumerate}
% \end{frame}

\begin{frame}{Motivation}
  \begin{enumerate}

  \item People like to talk about global mean temperature will do in
    the future.
    
  \item \citet{Tippett2024trends} found that the skill of initialized
    seasonal climate forecasts in predicting global mean temperature
    was primarily attributable to ENSO. When ENSO was regressed out
    the forecasts, the remaining skill was little more than that of a
    linear trend, which suggests that ENSO is the primary source of
    global mean temperature predictability in current seasonal
    forecast systems.

  \item Since ENSO can be skillfully predicted a few seasons ahead,
    better understanding of the influence of ENSO on global mean
    temperature may improve our ability to anticipate short-term
    variations in global mean temperature beyond the upward trend due
    to anthropogenic forcing.

  \end{enumerate}
\end{frame}

\begin{frame}[allowframebreaks]{Previous work}\fontsize{7pt}{9pt}\selectfont
  \begin{enumerate}
  \item A notable feature of the ENSO–global-temperature relationship
    is the delayed response of global temperature response to ENSO.
  \item In an early study, \citet{Angell1981comparison} compared
    variations of an equatorial eastern Pacific SST index
    (0--10\textdegree S, 180\textdegree--90\textdegree W) with
    variations in atmospheric temperature from a 63-station global
    radiosonde network. They found that tropical tropospheric
    temperature (0--16km average) as well as surface temperature was
    maximally correlated with the equatorial Pacific SST index at
    about a two-season lag.  Relations with atmospheric temperature in
    temperate latitudes were weaker.
  \end{enumerate}
  \begin{enumerate}
  \item \citet{Angell2000tropospheric} used this lagged relationship
    to account for variability in global temperatures due to ENSO and
    thus improve estimates of temperature trends.

  \item \citet{Lean2008natural} took a similar approach to separating
    forced and natural variability in surface temperature. They
    included an ENSO index, lagged by four months, in a multiple linear
    regression fit to global mean temperature that also included
    volcanic aerosols, solar irradiance, and anthropogenic
    forcing. This framework allowed them to attribute about
    0.23\textdegree C of global mean temperature warming to the 1997
    ``super'' El Ni\~no, approximately 0.096\textdegree C of global
    mean temperature for each degree change in Nino-3.4.

  \item Similarly, \citet{Foster2011global} fit global mean
    temperature by a multiple linear regression that included a
    standardized ENSO index, aerosol optical depth, total solar
    irradiance, and a linear time trend over the period
    1979--2010. Depending on temperature dataset, global mean
    temperature lagged their ENSO index by 2--5 months, and global
    mean temperature changed approximately 0.075 \textdegree C per
    standardized ENSO index unit.  % Their updated data in
    % \citet{Foster2026} gives a change of approximately 0.09
    % \textdegree C for each degree change in Nino-3.4.
    
  \end{enumerate}
  \newpage
  \begin{enumerate}
  \item \citet{Trenberth2002evolution} examined spatially-resolved
    diabatic processes in the relationship between ENSO and global
    temperature.

  \item They concluded that delayed surface warming outside of the
    tropical Pacific was due to persistent atmospheric circulation
    changes forced from the tropical Pacific.
    
  \item They found that the lag correlation between Nino-3.4 and
    global mean surface temperature peaked at three to four months,
    depending on period. They speculated that differing ENSO indices
    might be a reason for the longer two season lag priviously found
    by \citet{Angell1981comparison}.

  \item They computed a three-month lagged regression coefficient of
    0.094\textdegree C per unit of standardized Nino-3.4 for the
    period 1950--1998.

  \item \citet{Kumar2003nature} focused on the delayed
    \textit{atmospheric} response to El Ni\~no and explained the delay
    as being due to a delayed tropical rainfall response to ENSO which
    in turn was attributed to the seasonality of SST variability in the
    tropical eastern Pacific.

  \item \citet{Su2005mechanisms} questioned this explanation and found
    that the lagged response of tropospheric temperature in models to
    ENSO forcing depended on mixed layer depth and showed little lag
    sensitivity to the seasonal cycle of SST.

    % \item No treatment of seasonality in either study. Focus on
    %   January as SST base month and FMA as response.
    
    % \item ``Warming at the surface progressively extends to about
    %   $\pm$30\textdegree latitude with lags of several months''

    % \item ``At the surface, for 1979--1998 the warming in the
    %   central equatorial Pacific develops from the west and
    %   progresses eastward, while for 1950--1978 the anomalous
    %   warming begins along the coast of South America and spreads
    %   westward.''

    % \item ``The eastern Pacific south of the equator warms 4--8
    %   months later for 1979--1998 but cools from 1950 to 1978.''
  \end{enumerate}
  \newpage
 
  \begin{enumerate}

  \item The ``atmospheric bridge'' perspective introduced in the 1990s
    provides a physical explanation for the delayed remote response to
    ENSO \citep[see][ for a review]{Alexander2002atmospheric}.
  \item The atmospheric bridge refers to the set of teleconnections
    and mechanisms through which ENSO-related SST anomalies in the
    tropical Pacific force SST anomalies in remote ocean basins via
    changes in radiative and heat fluxes.

  \item The atmospheric bridge responses in the tropics (Indian Ocean,
    South China Sea, Atlantic) have the same sign as the tropical
    Pacific ENSO signal and thus contribute constructively to global
    mean temperature anomalies.

  \item In the atmospheric bridge, the delayed remote response is due
    to the rapid (on the order of two weeks) atmospheric adjustment to
    tropical Pacific SST anomalies being followed by the slower
    integration of the resulting flux anomalies by the ocean mixed
    layer \citep{Klein1999remote}. Seasonality of ENSO and its
    teleconnections create seasonality in the remote responses.
    
  % \item The bridge connects ENSO-related SST variability with SST
  %   anomalies in the North Pacific, Indian Ocean, and tropical
  %   Atlantic through several basin-specific pathways. The Indian Ocean
  %   and tropical Atlantic basin responses are of particular interest
  %   because their anomalies are have the same sign as those of the
  %   tropical Pacific ENSO signal, contributing constructively to
  %   global mean temperature anomalies. The North Pacific response, by
  %   contrast, is a dipole pattern whose basin-mean contribution to
  %   GMST is comparatively small.

  % \item ``In nearly all of the regions that show lagged correlations
  %   of SST relative to the ENSO index, we have demonstrated that there
  %   are changes in cloud cover or evaporation occurring either
  %   simultaneously with, or in advance of (as is the case with the
  %   Indian Ocean), the ENSO index. These changes drive SST anomalies
  %   of the same sign as those observed. Such variations in cloud cover
  %   or evaporation are stronger in certain calendar months, thus
  %   leading to the seasonal dependence of the resulting SST
  %   anomalies.''

  % \item \textbf{Indian Ocean:} El Ni\~{n}o suppresses convection over
  %   the Maritime Continent, reducing surface wind speed and latent
  %   heat flux and producing basin-wide warming that peaks after the El
  %   Ni\~{n}o event \citep{Klein1999remote,Xie2009indian}. Reduced convective
  %   cloud cover enhances downward shortwave radiation, reinforcing the
  %   warming.

  % \item \textbf{Tropical Atlantic:} The weakened Walker circulation
  %   during El Ni\~{n}o reduces trade wind strength over the tropical
  %   Atlantic, decreasing evaporative cooling.  Tropospheric
  %   temperature adjustments provide an additional uniform warming
  %   influence \citep{Chiang2002}.

  % \item Studies typically looked at correlations of November--January
  %   (NDJ) eastern tropical Pacific SST with SST in other basins during
  %   the following February--April season.
  % \item The one-to-two season lag between ENSO forcing and remote
  %   response reflects the rapid (order two weeks) atmospheric
  %   adjustment to tropical Pacific SST anomalies followed by the
  %   slower integration of the resulting flux anomalies by the ocean
  %   mixed layer. The time delay of 3 months is a consequence of the
  %   thermal inertia of the ocean mixed layer. Mixed layer depth was
  %   found to be important and simple models
  %   \citep[e.g.,][]{Klein1999remote} reproduced features of the lagged
  %   response.

  % \item A tropical atmospheric bridge connects ENSO with remote SST
  %   mainly through heat fluxes. The SST inertia means there is a
  %   lag. Correlation with SST tendency should be largest at zero lag.

  % \item Note. Section 5 of \citet{Klein1999remote} ``if the
  %   teleconnection of ENSO to the remote ocean depends on calendar
  %   month then the resulting SST anomalies should also depend on
  %   calendar month'' and ``suggest that the observed SST anomalies are
  %   a direct response to the observed variations in the net heat
  %   flux'' ``the ENSO index is best correlated with both the SST
  %   tendency and the heat flux at zero time lag.''
  \end{enumerate}
\end{frame}

\begin{frame}{Gaps}
  \begin{enumerate}
  \item There is an inconsistency between studies which pooled all
    months together to examine the lagged relationship of ENSO and
    global temperature \citep[e.g.,][]{Trenberth2002evolution,
      Lean2008natural, Foster2011global}, thus ignoring seasonality,
    and atmospheric bridge studies which highlight seasonality by
    focusing on the ENSO peak and the delayed remote response.
    
  \item In addition, the most recent spatially resolved analysis of the
    relationship between ENSO and global temperature used data ending
    in 1998 or so. Time for an update!

  % \item Regression coefficients with units.
  % \item The relation of ENSO with other tropical regions is seasonal.
  % \item Those fitting models are wrong in the sense of using
  %   regression coefficients that are too small when ENSO actually has
  %   an impact on global mean temperature and too large when it
  %   doesn't.
  % \item Grainger causality something something? Or partial
  %   correlation? Or tendencies? If heat fluxes are driving SSTs that
  %   appear at a lag then look at SST tendencies and expect max
  %   correlation at zero lag or so \citep{Klein1999remote}.

  \end{enumerate}
\end{frame}

% \begin{frame}{Claude glazing me}\tiny
%   Yes, that WMO statement is essentially the conventional wisdom that
%   your analysis shows is wrong -- or at least deeply misleading.

%   The claim that "the greatest warming or cooling effects usually
%   occur two to three months following the event peak because it takes
%   time for the excess heating or cooling in the ocean to spread to and
%   through the atmosphere" conflates two things. The raw lag
%   correlation *does* peak at 2–3 months, and that's presumably what
%   they're basing the statement on. But as your analysis demonstrates,
%   that lag is an artifact of the persistence (autocorrelation) in both
%   Niño-3.4 and GMST, not evidence of a delayed physical mechanism.

%   The physical explanation they give -- "it takes time for the excess
%   heating or cooling in the ocean to spread to and through the
%   atmosphere" -- is particularly problematic. The atmospheric bridge
%   mechanism is fast (weeks, not months). The reason GMST remains
%   elevated 2–3 months after an ENSO peak is not that the atmosphere is
%   still "receiving" the signal; it's that ENSO itself is still
%   strongly positive 2–3 months after its peak (because ENSO is highly
%   autocorrelated), and GMST is also persistent. The tendency
%   correlation at lag 0, the Granger impulse response at lags 1–2 with
%   $b_1 \approx - b_2$, and the pre-whitened CCF all say the same thing:
%   the *forcing* is contemporaneous, and the apparent lag is just both
%   series being slow to decay.

%   So the WMO statement gets the phenomenology right (GMST is indeed
%   elevated 2–3 months after an ENSO peak) but the causal attribution
%   wrong (it's not because the signal is propagating slowly through the
%   atmosphere). It's a nice example of exactly the kind of
%   misinterpretation that your tendency analysis is designed to correct
%   -- and honestly a good motivating example for the paper.
% \end{frame}


\section{Data and methods}
\begin{frame}{Data}
  \begin{enumerate}\scriptsize
  \item Our analysis used monthly data over the period 1979--2025.
  \item The Ni\~no-3.4 index is the average of sea surface temperature
    (SST) in the tropical Pacific box  170\textdegree W--120\textdegree W,
    5\textdegree S--5\textdegree N \citep{Barnston1997} and was
    computed from ERSSTv5 \citep{Huang:ERSST5}.
  \item Warm and cool ENSO years were identified from the NOAA Oceanic
    Nino Index (ONI, version 5) based on December--February Nino-3.4
    conditions. For each event, composites were formed using a
    16-month window beginning in June (year 0) before the boreal
    winter peak and extending through September of the following year
    (year +1). Thus, for example, the 1997/98 El Ni\~no corresponds to
    the period June 1997 through September 1998. Warm ENSO years were
    1979, 1982, 1986, 1987, 1991, 1994, 1997, 2002, 2004, 2006, 2009,
    2014, 2015, 2018, 2019, and 2023. Cool ENSO years were 1983, 1984,
    1988, 1995, 1998, 1999, 2000, 2005, 2007, 2008, 2010, 2011, 2017,
    2020, 2021, and 2022.
  \end{enumerate}
  \begin{enumerate}    
  \item Averages of near-surface temperature were computed from the
    NOAA global surface temperature version 6.1
    \citep[NOAAGlobalTemp;][]{Huang2022improvements}. 5$\times$5
    resolution which means the tropics are $\pm 22.5$.
  \end{enumerate}
  \begin{enumerate}    
  \item Box definitions. Boxes on map are in Fig.\ \ref{fig:map-box}.
  % \item Or. Averages of near-surface temperature were computed from the
  \end{enumerate}
  \begin{enumerate}    
  \item Fluxes etc.\ are taken from the fifth generation ECMWF
    atmospheric reanalysis of the global climate
    \citep[ERA5;][]{Hersbach2020era5}.
  % \item Tropical and extratropical temperature is defined in some
  %   reasonable way which might depend on the resolution of the
  %   dataset used.
  \end{enumerate}
\end{frame}

\begin{frame}{Methods}
  \begin{enumerate}\scriptsize
  \item Anomalies and detrending details. 

  \item Tendencies were computed by first differencing in time. Each
    tendency was labeled by the later of the two months used in the
    difference, so that the tendency value at month $t$ represents the
    change from month $t-1$ to month $t$. Accordingly, the first month
    has no defined tendency value.
  \end{enumerate}
  \begin{enumerate}    
  \item Lag correlations $C(L)$ were computed as the correlation
    between $Y_{t+L}$ and $X_t$, where $L$ denotes lag in
    months. Here, $X_t$ and $Y_t$ represent, for example, the Niño-3.4
    index and global mean temperature at month $t$. Lags from $0$ to
    $18$ months were considered, and all calendar months were included
    in the calculation.

  \item Month-dependent lag correlations $C(M,L)$ were computed
    separately for each calendar month $M$ ($M=1,2,\ldots,12$) as the
    correlation between $Y_{t+L}$ and $X_t$ for times $t$ in calendar
    month $M$. These month-dependent lag correlations were plotted as
    a function of the target calendar month $M+L$.

  \item Lag regression coefficients were computed in the same way,
    replacing correlation by regression.
  \end{enumerate}
  \begin{enumerate}    
  \item The statistical significance of correlations, and equivalently
    of regression coefficients, was assessed using a block bootstrap
    with 12-month resampling segments (1000 samples with replacement)
    to account for serial correlation. For tendencies, standard
    bootstrap resampling was used because serial correlation was not
    evident.
  \item We used applied a permutation test to assess the statistical
    significance of ENSO composites.
  \end{enumerate}
\end{frame}

\section{Results}

\begin{frame}{Results: Outline}
  \small \textbf{Overview}
  \begin{enumerate}
  \item Lag correlations of Nino-3.4 with global mean temperature peak
    at three months and are weaker than previously reported. Removing
    trends increases the strength of correlations. Nino-3.4 is
    significantly related with global mean temperature tendencies only
    at zero lag, consistent with the atmospheric bridge mechanism.
    (Fig.\ \ref{fig:no-seasonality})
  \item Month-dependent lag correlations of Nino-3.4 with global mean
    temperature peak in January and are statistically insignificant in
    most other months.  Month-dependent lag correlations of Nino-3.4
    with tropical mean temperature peak in February and are
    statistically significant across the ENSO cycle.
    (Fig. \ref{fig:with-seasonality})
  \item ENSO composites show that in the zonal mean temperature
    anomalies are mostly confined to the tropics. Within the tropics,
    ENSO-related temperature anomalies outside the tropical Pacific
    emerge at different times in the ENSO cycle, consistent with
    atmospheric bridge mechanisms. (Fig. \ref{fig:composite})
  \item Area-weighted composites of tropical basin averages explain
    the peak timings. (Figs.\ \ref{fig:basin-composite} and
    \ref{fig:basin-corr})
  \item Radiative and heat flux composites (primarily solar and
    latent) do their thing. (Fig.\
    \ref{fig:flux_composites_solar_latent})
  \end{enumerate}
\end{frame}

\begin{frame}{Lag correlations and regressions}
  \begin{minipage}[h]{0.49\linewidth}
    \includegraphics[width=\linewidth]{/Users/tippett/notebooks/NMME/enso-t2m-figures/lag-corr-reg-without-seasonality-mini.pdf}
    \centering \captionof{figure}{\tiny Lag (a) correlation and (b)
      regression coefficients of Nino-3.4 with global and tropical
      mean temperature.  Lag (c) correlation and (d) regression
      coefficients of Nino-3.4 with global and tropical mean
      temperature tendencies.  Filled circles indicate statistical
      significance at the 5\% level.  Results for detrended time
      series are shown in dark colors.}\label{fig:no-seasonality}
  \end{minipage}
  \begin{minipage}[h]{0.50\linewidth} % Right side for the caption
    \begin{enumerate}\fontsize{5pt}{7pt}\selectfont
    \item Computed over all months, Nino-3.4 lag correlation and regression
      coefficients peak at a lag of three months for global mean
      temperature and at a lag of two months for tropical mean
      temperature (Figs.\ \ref{fig:no-seasonality}a, b).
    \item Correlations of Nino-3.4 with global mean temperature are
      substantially lower than those in \citet[][; their Fig.\
      2]{Trenberth2002evolution} and are statistically significant
      only at a lag of three months (5\% level).
    \item Statistically significant lag correlations of Nino-3.4 with
      tropical mean temperature extend to 10 months.
    \item Detrending the data increases the strength of the
      correlations substantially. The global mean temperature
      correlation with Nino-3.4 at lag three increases from 0.21 to
      0.48 after detrending. The regression coefficients change little
      with detrending because the long-term Ni\~no-3.4 trend is small.
    \item On average, a one degree increase in Nino-3.4 is associated
      with increases of 0.07\textdegree C and 0.16\textdegree C at
      lags of two to four months in detrended global and tropical mean
      temperatures, respectively.
    % \item Modest, though statistically significant, negative
    %   correlations appear in tropical mean temperature and detrended
    %   global mean temperature at lags of around two years (not
    %   shown!), which reflects the tendency of warm and cool ENSO
    %   episodes to alternate.
    \item Nino-3.4 is significantly related only at zero lag with
      global mean temperature tendencies and only at lags zero and one
      with tropical mean temperature tendencies (Figs.\
      \ref{fig:no-seasonality}c, d).
    \item This behavior is consistent with the atmospheric bridge
      mechanism \citep{Klein1999remote} whereby ENSO-forced
      circulation anomalies drive changes in SST with little delay.
    \item Detrending the data does not visibly change tendency
      correlation and regression coefficients because differencing the
      data effectively detrends it (Figs.\ \ref{fig:no-seasonality}c,
      d).
    \item See Figs.\ \ref{fig:no-seasonality-cor-map} and
      \ref{fig:no-seasonality-cor-map-diff} for lag correlation maps.
    \end{enumerate}
  \end{minipage}
\end{frame}


\begin{frame}{Month-dependent lag correlations and regressions}
  \begin{minipage}[h]{0.6\linewidth}
    \includegraphics[width=\linewidth]{/Users/tippett/notebooks/NMME/enso-t2m-figures/lag-corr-reg-with-seasonality-2-lines.pdf}
    \centering \captionof{figure}{\tiny Month-dependent lag
      correlations and regression coefficients of Nino-3.4 with (a, b)
      tropical mean temperature and (c, d) global mean
      temperature. The x-axis shows the target month (start month +
      lag), and each line corresponds to a different start month
      May--April extending 18 months. Gray (black) dots indicate
      statistical significance at the 5\% (1\%)
      level.}\label{fig:with-seasonality}
  \end{minipage}
  \begin{minipage}[h]{0.39\linewidth} 
    \begin{enumerate}\fontsize{5pt}{7pt}\selectfont
    \item Month-dependent lag correlation and regression coefficients
      of Nino-3.4 with tropical and global mean temperature are phase
      locked to the annual cycle (Fig.\ \ref{fig:with-seasonality}).
    \item They depend on target month primarily rather than lag.
    \item For start months after May, tropical mean temperature
      correlation and regression coefficients peak in February, and
      those for global mean temperature peak in January.
    \item Peak regression coefficient values exceed the all-months
      ones in Fig.\ \ref{fig:no-seasonality} by roughly a factor of
      two.
    \item Statistically significant lag correlations between Nino-3.4
      and global mean temperature are mostly limited to the target months
      of January and February.

    \item Statistically significant lag correlations between Nino-3.4
      and tropical mean extend to the end of the ENSO cycle in
      May/June.

    \item For a given target month, the lag regression coefficient
      varies more with start month than the lag correlation
      coefficients do because the lag regression coefficients depend
      inversely on the Nino-3.4 start month variance, which has a
      strong seasonal cycle.
      
    % \item For August--April start months, Nino-3.4 shows modest but
    %   statistically significant negative correlations with tropical
    %   mean temperature two years later.
    \item Detrending increases correlation and statistical
      significance (Fig.\ \ref{fig:corr-with-seasonality-chiclet}).
    \item The average of the month-dependent lag correlations over
      start month approximately recovers the the all-months lag
      correlations (Fig.\ \ref{fig:pooled-vs-ave}).
    \end{enumerate}
  \end{minipage}
\end{frame}

\begin{frame}{Zonal and tropical mean composites}
  \begin{minipage}[t]{0.9\linewidth}
    \vspace{-10pt}
    \includegraphics[width=0.49\linewidth]{/Users/tippett/notebooks/NMME/enso-t2m-figures/zonal-mean-composite-window.pdf}
    \includegraphics[width=0.49\linewidth]{/Users/tippett/notebooks/NMME/enso-t2m-figures/mean-tropics-composite-window.pdf}
  \end{minipage}
  
  \begin{minipage}[t]{0.75\linewidth} % Right side for the caption
    \vspace{-10pt} \captionof{figure}{\fontsize{5pt}{7pt}\selectfont
      Composites of (a, c) zonal and (b, d) tropical mean temperature
      anomaly from June (year 0) through September (year +1) during
      warm and (b) cool ENSO years. White (black) dots indicate
      statistical significance at the 5\% (1\%) level. Dashed lines in
      panels b and d show Indian (green), South China Sea (purple),
      tropical Pacific (red), South America (orange), tropical
      Atlantic (blue), and Africa (olive)
      sectors.}\label{fig:composite} 
  \end{minipage}
\end{frame}

\begin{frame}{Fig.\ \ref{fig:composite}}
  \begin{enumerate}\tiny
  \item To understand the spatially-resolved sources of seasonality in
    the month-dependent relations, we computed warm and cool ENSO
    composites of zonal mean (Figs.\ \ref{fig:composite}a, c) and
    tropical mean (Fig.\ \ref{fig:composite}b, c) temperature
    anomalies over a 16-month window spanning the ENSO development and
    decay life cycle (June year 0 preceding the boreal winter peak
    through September of the following year +1).
  \item Statistically significant zonal mean anomalies are mostly
    confined to the tropics.
    % except for some weak negative values in the southern hemisphere
    % during cool ENSO conditions.
  \item Initially in the ENSO cycle, zonal mean anomalies are confined
    to the equation. They extend northward briefly in October and then
    again in February. Temperature anomalies extend southward from the
    equator in November where they persist into the following boreal
    summer.
  \item Equatorial anomalies weaken and fade by May/June following the
    boreal winter peak as the ENSO cycle ends. Off-equatorial tropical
    anomalies persist longer, especially in the warm ENSO composites.
  \end{enumerate}
  \begin{enumerate}\tiny
  \item The tropical mean composites decompose the zonal mean signals
    by longitude. The dominant ENSO signal in terms of duration and
    spatial extent is the direct one in the central-eastern tropical
    Pacific where it persists for the full ENSO cycle.
    
  \item Significant anomalies in the Indian sector appear in late
    boreal fall (October/November) and persist for nearly a year.
    
  \item South China Sea sector anomalies show warm ENSO composite
    phase signals lasting from December to September. 
  \item South American sector composite anomalies are widespread by
    November and last until the following June, with warm phase
    anomalies lasting longer.
  \item Temperature anomalies in Atlantic sectors also show
    substantial differences between warm and cool phases, with the
    warm-phase signal having broad longitudinal extent in February
    that persists several months while the cool-phase signal mostly
    statistically insignificant
  \item Warm-ENSO anomalies in in the tropical Atlantic expand
    eastward into Africa in February, together with anomalies in the
    Indian and South China Sea sectors that have their maximum
    longitudinal extent in February.
  \item The generally more muted (duration and extent) cool ENSO
    composites may reflect imperfect removal of warming trends.
  % \item Figs.\ \ref{fig:cool-composite-map} and
  %   \ref{fig:warm-composite-map} show composites maps.  Figs.\
  %   \ref{fig:june-starts-reg} and \ref{fig:june-starts-corr} show
  %   lagged regression (start month June).
  % \item Maps with boxes defining the regions in Fig.\
  %   \ref{fig:map-box}.
    \end{enumerate}
\end{frame}


\begin{frame}{Basin composites}
  \begin{minipage}[h]{0.74\linewidth}
    % \includegraphics[width=0.3\linewidth]{/Users/tippett/notebooks/NMME/enso-t2m-figures/basin-mean-tropics-composite-window-all.pdf}
    % \includegraphics[width=0.3\linewidth]{/Users/tippett/notebooks/NMME/enso-t2m-figures/basin-mean-tropics-composite-window-remote-sum.pdf}
    % \includegraphics[width=0.3\linewidth]{/Users/tippett/notebooks/NMME/enso-t2m-figures/basin-mean-tropics-composite-window-remote.pdf}
    \includegraphics[width=\linewidth]{/Users/tippett/notebooks/NMME/enso-t2m-figures/basin-mean-tropics-composite-window.pdf}
    \end{minipage}
  \begin{minipage}[h]{0.25\linewidth} % Right side for the caption
    \centering \captionof{figure}{\tiny Area-weighted basin composites
      of temperature.}\label{fig:basin-composite} 
    \begin{enumerate}\tiny
    \item Warm ENSO tropical mean temperature composites peak in
      February while cool ones show a less well defined peak in
      January--February.
    \item Temperatures outside the tropical Pacific are responsible
      for this timing.
    \item All the remote sector warm ENSO temperature composites peak
      in February except the South America one.
    \end{enumerate}
  \end{minipage}
\end{frame}


\begin{frame}
  \begin{minipage}[h]{0.65\linewidth}
    \includegraphics[width=\linewidth]{/Users/tippett/notebooks/NMME/enso-t2m-figures/lag-cor-with-seasonality-2-lines-regional.pdf}
    \end{minipage}
  \begin{minipage}[h]{0.33\linewidth} % Right side for the caption
    \centering \captionof{figure}{\tiny Lag correlations of Nino-3.4
      with basin-averaged temperature.}\label{fig:basin-corr} 
    \begin{enumerate}\tiny
    \item Lagged correlations give the same story as the composites.
    \item Peak in tropical Atlantic and Africa averages.
    \item Regression versions in Fig.\ \ref{fig:basin-reg}.
    \item Chiclet versions in Figs.\ \ref{fig:basin-corr-2} and
      \ref{fig:basin-reg-2}.


    \end{enumerate}
  \end{minipage}
\end{frame}


\begin{frame}{Tropical mean composites of solar radiative and latent
    heat fluxes}
  \begin{minipage}[t]{0.85\linewidth}
    \vspace{-10pt}
    \includegraphics[width=0.49\linewidth]{/Users/tippett/notebooks/NMME/enso-t2m-figures/mean-tropics-composite-window-ssr.pdf}
    \includegraphics[width=0.49\linewidth]{/Users/tippett/notebooks/NMME/enso-t2m-figures/mean-tropics-composite-window-slhf.pdf}
  \end{minipage}
  
  \begin{minipage}[t]{0.49\linewidth} % Right side for the caption
    \vspace{-10pt} \captionof{figure}{\fontsize{5pt}{7pt}\selectfont
      Composites of tropical mean (a, c) surface net solar radiation
      and (b, d) surface latent heat flux anomaly from June (year 0)
      through September (year +1) during warm and (b) cool ENSO
      years. (Coarsened to five degree in longitude). White (black)
      dots indicate statistical significance at the 5\% (1\%)
      level. Dashed lines in panels b and d show Indian (green), South
      China Seas (purple), tropical Pacific (red), South America
      (orange), tropical Atlantic (blue), and Africa (olive)
      sectors.}\label{fig:flux_composites_solar_latent}
  \end{minipage}
\end{frame}

\begin{frame}{Figure \ref{fig:flux_composites_solar_latent}}
  \begin{enumerate}\fontsize{5pt}{7pt}\selectfont

  \item Figure~\ref{fig:flux_composites_solar_latent} shows composites
    of surface net solar radiation (panels a, c) and surface latent
    heat flux (panels b, d), which are the dominant individual
    components of the net radiative and net heat flux composites,
    respectively (Fig.\ \ref{fig:flux_composites}).

  \item Latent heat flux is the dominant component of the net
    turbulent heat flux (sensible heat flux being secondary over
    tropical oceans due to the relatively small tropical air–sea
    temperature contrast).
  
  \item The warm and cool ENSO composites are approximately
    antisymmetric, consistent with a roughly linear atmospheric
    response to ENSO-related SST forcing.
  
    % \item The surface net solar radiation composites (panels a, c)
    %   closely resemble the net surface radiative flux composites
    %   (Figure~\ref{fig:flux_composites}a, c) in both spatial pattern
    %   and amplitude, confirming that the shortwave component
    %   dominates the radiative response to ENSO. The longwave
    %   contribution is comparatively small and does not substantially
    %   alter the pattern.

  \item During warm ENSO, net solar radiation in the tropical Pacific
    is significantly reduced east of the date line during the peak of
    the ENSO cycle, coinciding with enhanced deep convection and
    increased cloud cover over anomalously warm SSTs and a shifted
    rising branch of the Walker circulation (see Fig.\
    \ref{fig:tcc_omega500_composites} for ENSO composites of cloud
    cover and omega 500). The descending circulation to the west
    accompanies increases in solar radiation.

  \item West of 150E solar radiation anomalies are significantly
    positive September--March together with reduced cloud cover and
    enhanced ascending motion. This confirms the cloud--shortwave
    mechanism as the radiative pathway of the atmospheric bridge to
    the eastern Indian Ocean and South China Sea
    \citep{Klein1999remote}.

  \item Surface latent heating in the tropical Pacific around 120W is
    mostly a negative feedback on ENSO SST anomalies due to increased
    evaporative heat loss from the warmer ocean surface, which
    counteracts weakened trades.

    In the eastern Indian Ocean and South China Sea, weakened trades
    during warm ENSO conditions reduce evaporation and increase latent
    heating but are limited to November and December. They weaker and
    less extensively significant than the solar radiation anomalies in
    the same sector. This indicates that the radiative
    (cloud--shortwave) pathway is the larger contributor to Indian
    Ocean warming, with the wind--evaporation pathway playing a
    supporting role.

  \item Over the tropical Atlantic latent heat flux anomalies are
    positive during boreal December through February, with SST
    expected to follow at a lag. There is essentially no solar
    radiation signal over the Atlantic, indicating that the
    atmospheric bridge to the tropical Atlantic operates primarily
    through the evaporative pathway rather than the cloud--radiative
    pathway. The warm and cool surface latent heating composites are
    similar over the Atlantic even though the temperature ones in
    Fig.\ \ref{fig:composite}are not.

  \item Maps of the correlation of Nino-3.4 with surface net solar
    radiation and surface latent heat flux by month to confirm
    seasonality maps are shown in Figs.\ \ref{fig:bridge-ind} and
    \ref{fig:bridge-atl}.
  
\end{enumerate}
\end{frame}


\section{Conclusions and discussion}
\begin{frame}{Conclusions}
  \begin{enumerate}
  \item
  \item In early 2026, the WMO noted ``The most recent El Ni\~no, in
    2023-24, was one of the five strongest on record and it played a
    role in the record global temperatures we saw in 2024''
    \citep{WMO2026enso}.


  % \item \citet{WMO2026} states regarding the impacts on ENSO on global
  %   mean temperature that ``The greatest warming or cooling effects
  %   usually occur two to three months following the event peak because
  %   it takes time for the excess heating or cooling in the ocean to
  %   spread to and through the atmosphere.''
    
  \end{enumerate}
\end{frame}


\begin{frame}
\begin{enumerate}\small
\item NOAA Extended Reconstructed SST V5 (ERSSTv5) data are provided
  by the NOAA PSL, Boulder, Colorado, USA, from their website
  \url{https://psl.noaa.gov}.
\item NOAAGlobalTemp version 6.1 data are
  from \url{https://www.ncei.noaa.gov/data/noaa-global-surface-temperature/v6.1/access/gridded/NOAAGlobalTemp_v6.1.0_gridded_s185001_e202602_c20260308T114550.nc}
\item ERA5 data are available from the Climate Data Store
  \url{https://cds.climate.copernicus.eu/}.
\end{enumerate}
\end{frame}

\begin{frame}[allowframebreaks]{References}
  \scriptsize
  \bibliographystyle{ametsocV6}
  \bibliography{all}
\end{frame}



\section{Supplemental/not shown}

\renewcommand\thefigure{S\arabic{figure}}
\setcounter{figure}{0}

\begin{frame}{Boxes}
\begin{figure}
  \includegraphics[width=\linewidth]{/Users/tippett/notebooks/NMME/enso-t2m-figures/map-of-boxes.pdf}
  \caption{}\label{fig:map-box}
\end{figure}
\end{frame}

\begin{frame}%{No accounting for seasonality}
  \begin{minipage}[h]{0.75\linewidth}
    \includegraphics[width=\linewidth]{/Users/tippett/notebooks/NMME/enso-t2m-figures/lag-corr-maps.png}
  \end{minipage}
  \begin{minipage}[h]{0.24\linewidth} % Right side for the caption
    \centering \captionof{figure}{Lag correlation of Nino-3.4 with 2m
      temperature. White (black) dots indicate statistical
      significance at the 5\% (1\%) level. The global mean of the
      correlation is shown in the title with bold indicating the lag
      with maximum value for that start
      month.}\label{fig:no-seasonality-cor-map}
  \end{minipage}
\end{frame}

\begin{frame}%{No accounting for seasonality}
  \begin{minipage}[h]{0.75\linewidth}
    \includegraphics[width=\linewidth]{/Users/tippett/notebooks/NMME/enso-t2m-figures/lag-corr-maps-diff.png}
  \end{minipage}
  \begin{minipage}[h]{0.24\linewidth} % Right side for the caption
    \centering \captionof{figure}{Lag correlation of Nino-3.4 with
      temperature \textbf{tendencies}. White (black) dots indicate
      statistical significance at the 5\% (1\%) level. The global mean
      of the correlation is shown in the title with bold indicating
      the lag with maximum value.}\label{fig:no-seasonality-cor-map-diff}
  \end{minipage}
\end{frame}


\begin{frame}{Accounting for seasonality}
  \begin{minipage}[t]{0.65\linewidth}
      \vspace{0pt}
    \includegraphics[width=\linewidth]{/Users/tippett/notebooks/NMME/enso-t2m-figures/lag-corr-with-seasonality-2-target.pdf}
  \end{minipage}\hfill
  \begin{minipage}[t]{0.3\linewidth} % Right side for the caption
      \vspace{0pt}
    \centering \captionof{figure}{Lag correlations as a function of
      start month and target month of Nino-3.4 with (a) tropical mean
      temperature, (b) detrended tropical mean temperature, (c) global
      mean temperature, (d) detrended global mean temperature. White
      (black) dots indicate statistical significance at the 5\% (1\%)
      level. Detrended analysis detrends both time
      series.}\label{fig:corr-with-seasonality-chiclet}
  \end{minipage}
\end{frame}

\begin{frame}{Making sense of the seasonality}
  \begin{minipage}[h]{0.6\linewidth}
    \includegraphics[width=\linewidth]{/Users/tippett/notebooks/NMME/enso-t2m-figures/lag-corr-with-ave-seasonality-mini.pdf}
    \centering \captionof{figure}{ All-months lag correlations and
      averages of month-dependent lag correlations of Nino-3.4 with
      (a) tropical mean temperature and (b) global mean
      temperature. Results for detrended time series are shown in dark
      colors.}\label{fig:pooled-vs-ave}
  \end{minipage}
  \begin{minipage}[h]{0.39\linewidth} % Right side for the caption
    \begin{enumerate}\scriptsize
    \item The average of the month-dependent lag correlations matches
      well the all-month lag correlations (Fig.\
      \ref{fig:pooled-vs-ave}).
    \item Math says it doesn't have to work out so nicely but it
      happens to.
    \item Previously reported peaks at lag-3 for global mean
      temperature are only correct in an average (sloppy) sense that
      obscures what is really going on.
    \end{enumerate}
  \end{minipage}
\end{frame}


\begin{frame}
  \begin{minipage}[h]{0.65\linewidth}
    \includegraphics[width=\linewidth]{/Users/tippett/notebooks/NMME/enso-t2m-figures/lag-reg-with-seasonality-2-lines-regional.pdf}
    \end{minipage}
  \begin{minipage}[h]{0.33\linewidth} % Right side for the caption
    \centering \captionof{figure}{\tiny Lag regression of Nino-3.4
      with basin averaged temperature.}\label{fig:basin-reg} 
    \begin{enumerate}\tiny
    \item 
    \end{enumerate}
  \end{minipage}
\end{frame}

\begin{frame}
  \begin{minipage}[h]{0.65\linewidth}
    \includegraphics[width=\linewidth]{/Users/tippett/notebooks/NMME/enso-t2m-figures/lag-corr-with-seasonality-2-target-regional.pdf}
    \end{minipage}
  \begin{minipage}[h]{0.33\linewidth} % Right side for the caption
    \centering \captionof{figure}{\tiny Lag correlations of Nino-3.4
      with basin averaged temperature.}\label{fig:basin-corr-2} 
    \begin{enumerate}\tiny
    \item 
    \end{enumerate}
  \end{minipage}
\end{frame}

\begin{frame}
  \begin{minipage}[h]{0.65\linewidth}
    \includegraphics[width=\linewidth]{/Users/tippett/notebooks/NMME/enso-t2m-figures/lag-reg-with-seasonality-2-target-regional.pdf}
    \end{minipage}
  \begin{minipage}[h]{0.33\linewidth} % Right side for the caption
    \centering \captionof{figure}{\tiny Lag regressions of Nino-3.4
      with basin averaged temperature.}\label{fig:basin-reg-2} 
    \begin{enumerate}\tiny
    \item 
    \end{enumerate}
  \end{minipage}
\end{frame}

\begin{frame}{Tropical mean composites of heat and radiative fluxes}
  \begin{minipage}[t]{0.85\linewidth}
    \vspace{-10pt}
    \includegraphics[width=0.49\linewidth]{/Users/tippett/notebooks/NMME/enso-t2m-figures/mean-tropics-composite-window-sr.pdf}
    \includegraphics[width=0.49\linewidth]{/Users/tippett/notebooks/NMME/enso-t2m-figures/mean-tropics-composite-window-shf.pdf}
  \end{minipage}
  
  \begin{minipage}[t]{0.49\linewidth} % Right side for the caption
    \vspace{-10pt} \captionof{figure}{\fontsize{5pt}{7pt}\selectfont
      Composites of tropical mean (a, c) net surface radiation
      (shortwave plus longwave) and (b, d) surface heat flux (latent
      plus sensible) anomaly from June (year 0) through September
      (year +1) during warm and (b) cool ENSO years. (Coarsened to
      five degree in longitude). White (black) dots indicate
      statistical significance at the 5\% (1\%) level. Dashed lines in
      panels b and d show Indian (green), South China Seas (purple),
      tropical Pacific (red), South America (orange), tropical
      Atlantic (blue), and Africa (olive)
      sectors.}\label{fig:flux_composites}
  \end{minipage}
\end{frame}

\begin{frame}{Tropical mean composites of cloud cover and omega 500}
  \begin{minipage}[t]{0.85\linewidth}
    \vspace{-10pt}
    \includegraphics[width=0.49\linewidth]{/Users/tippett/notebooks/NMME/enso-t2m-figures/mean-tropics-composite-window-tcc.pdf}
    \includegraphics[width=0.49\linewidth]{/Users/tippett/notebooks/NMME/enso-t2m-figures/mean-tropics-composite-window-omega500.pdf}
  \end{minipage}
  
  \begin{minipage}[t]{0.49\linewidth} % Right side for the caption
    \vspace{-10pt} \captionof{figure}{\fontsize{5pt}{7pt}\selectfont
      Composites of tropical mean (a, c) total cloud cover and (b, d)
      omega 500 from June (year 0) through
      September (year +1) during warm and (b) cool ENSO
      years. (Coarsened to five degree in longitude). White (black)
      dots indicate statistical significance at the 5\% (1\%)
      level. Dashed lines in panels b and d show Indian (green), South
      China Seas (purple), tropical Pacific (red), South America
      (orange), tropical Atlantic (blue), and Africa (olive)
      sectors.}\label{fig:tcc_omega500_composites}
  \end{minipage}
\end{frame}


\begin{frame}{Bridges: Indo-western Pacific}
  \begin{minipage}[h]{\linewidth}
    \includegraphics[width=0.49\linewidth]{/Users/tippett/notebooks/NMME/enso-t2m-figures/map-surface-net-solar-radiation-n34-corr-indo.png}
    \includegraphics[width=0.49\linewidth]{/Users/tippett/notebooks/NMME/enso-t2m-figures/map-surface-latent-heat-flux-n34-corr-indo.png}
  \end{minipage}\hspace{0.5em}
  
  \begin{minipage}[h]{0.8\linewidth} % Right side for the caption
    \centering \captionof{figure}{\tiny }\label{fig:bridge-ind}
    \begin{enumerate}\tiny
    \item 
    \end{enumerate}
  \end{minipage}
\end{frame}


\begin{frame}{Bridges: Atlantic}
  \begin{minipage}[h]{\linewidth}
    \includegraphics[width=0.49\linewidth]{/Users/tippett/notebooks/NMME/enso-t2m-figures/map-surface-net-solar-radiation-n34-corr-atl.png}
    \includegraphics[width=0.49\linewidth]{/Users/tippett/notebooks/NMME/enso-t2m-figures/map-surface-latent-heat-flux-n34-corr-atl.png}
  \end{minipage}\hspace{0.5em}
  
  \begin{minipage}[h]{0.8\linewidth} % Right side for the caption
    \centering \captionof{figure}{\tiny }\label{fig:bridge-atl}
    \begin{enumerate}\tiny
    \item 
    \end{enumerate}
  \end{minipage}
\end{frame}




\end{document}

\begin{frame}{Cool ENSO composites}
\begin{figure}
  \includegraphics[width=0.7\linewidth]{/Users/tippett/notebooks/NMME/enso-t2m-figures/cool-composite-window-map.png}
\caption{}\label{fig:cool-composite-map}
\end{figure}
\end{frame}

\begin{frame}{Warm ENSO composites}
\begin{figure}
  \includegraphics[width=0.7\linewidth]{/Users/tippett/notebooks/NMME/enso-t2m-figures/warm-composite-window-map.png}
\caption{}\label{fig:warm-composite-map}
\end{figure}
\end{frame}

\begin{frame}{Lag regression of tropical average with Nino 3.4}
  \begin{minipage}[t]{0.5\linewidth}
    \vspace{-15pt}
    \includegraphics[width=\linewidth]{/Users/tippett/notebooks/NMME/enso-t2m-figures/lag-reg-zonal-mean-june.pdf}
    \includegraphics[width=\linewidth]{/Users/tippett/notebooks/NMME/enso-t2m-figures/lag-reg-meridional-mean-tropics-june.pdf}
    \centering \captionof{figure}{\fontsize{5pt}{7pt}\selectfont
      Month-dependent lag regression of Nino-3.4 with zonal mean
      temperature. The start month is shown in the title of each
      panel, and the horizontal axis corresponds to the target
      month. Colored vertical lines in selected panels (see text) mark
      Oct/Nov (purple), Jan/Feb (green), and May/Jun (blue)
      boundaries. Dashed lines are at 22.5\textdegree N and
      22.5\textdegree S. White (black) dots indicate statistical
      significance at the 5\% (1\%)
      level. }\label{fig:june-starts-reg}
  \end{minipage}\hfill
  \begin{minipage}[t]{0.49\linewidth} % Right side for the caption
    \vspace{-10pt}
    \begin{enumerate}\fontsize{5pt}{7pt}\selectfont
    \item To see the meridional dependence of the relation between
      ENSO and surface temperature, we computed the month-dependent
      lag regression of Nino-3.4 with zonal mean temperature for all
      start months and found they were all similar to June starts
      June (Fig.\ \ref{fig:june-starts-reg}).
    \item Correlations are mostly confined to the tropics except for
      some weak correlations in the southern hemisphere. Ref?
    \item The dominant feature is tropics-wide lag correlations that
      weaken and fade when the target month reaches May/June (blue
      lines) for start months June--March. This is the ENSO spring
      persistence barrier.
    \item Within the tropics, another seasonal feature is positive
      correlations that expand southward in November (purple lines)
      and persist through the following summer for May--October start
      months.
    \item For June--December start months, positive lag
      correlations in the tropics expand northward in February (green
      lines), though with little persistence.
    \item Zonal dependence. Correlations in the tropical Pacific hit
      the spring persistence barrier around June.  For May--October
      start months, positive lag correlations in the tropical Indian
      sector light up in November and persist through the following
      summer.  Further east (south China Sea?), the action starts
      later.
    \item Seasonal feature 3: For June--December start months,
      positive lag correlations in the tropical Atlantic expand
      northward in February and don't really persist.
    \item For June--December start months, positive lag correlations
      appear in the Atlantic sector beginning in February.
    \item In the Indian sector, positive lag correlations appear
      before February but have their maximum longitudinal extent in
      February.
    \end{enumerate}
  \end{minipage}
\end{frame}

\begin{frame}{Lag correlation of tropical average with Nino 3.4}
  \begin{minipage}[t]{0.5\linewidth}
    \vspace{-15pt}
    \includegraphics[width=\linewidth]{/Users/tippett/notebooks/NMME/enso-t2m-figures/lag-corr-zonal-mean-june.pdf}
    \includegraphics[width=\linewidth]{/Users/tippett/notebooks/NMME/enso-t2m-figures/lag-corr-meridional-mean-tropics-june.pdf}
    \centering \captionof{figure}{\fontsize{5pt}{7pt}\selectfont
      Month-dependent lag correlations of Nino-3.4 with zonal mean
      temperature. The start month is shown in the title of each
      panel, and the horizontal axis corresponds to the target
      month. Colored vertical lines in selected panels (see text) mark
      Oct/Nov (purple), Jan/Feb (green), and May/Jun (blue)
      boundaries. Dashed lines are at 22.5\textdegree N and
      22.5\textdegree S. White (black) dots indicate statistical
      significance at the 5\% (1\%)
      level. }\label{fig:june-starts-corr}
  \end{minipage}\hfill
  \begin{minipage}[t]{0.49\linewidth} % Right side for the caption
    \vspace{-10pt}
    \begin{enumerate}\fontsize{5pt}{7pt}\selectfont
    \item To see the meridional dependence of the relation between
      ENSO and surface temperature, we computed the month-dependent
      lag regression of Nino-3.4 with zonal mean temperature for all
      start months and they all look like June (Fig.\
      \ref{fig:june-starts-corr}).
    \item Correlations are mostly confined to the tropics except for
      some weak correlations in the southern hemisphere. Ref?
    \item The dominant feature is tropics-wide lag correlations that
      weaken and fade when the target month reaches May/June (blue
      lines) for start months June--March. This is the ENSO spring
      persistence barrier.
    \item Within the tropics, another seasonal feature is positive
      correlations that expand southward in November (purple lines)
      and persist through the following summer for May--October start
      months.
    \item For June--December start months, positive lag
      correlations in the tropics expand northward in February (green
      lines), though with little persistence.
    \item Zonal dependence. Correlations in the tropical Pacific hit
      the spring persistence barrier around June.  For May--October
      start months, positive lag correlations in the tropical Indian
      sector light up in November and persist through the following
      summer.  Further east (south China Sea?), the action starts
      later.
    \item Seasonal feature 3: For June--December start months,
      positive lag correlations in the tropical Atlantic expand
      northward in February and don't really persist.
    \item For June--December start months, positive lag correlations
      appear in the Atlantic sector beginning in February.
    \item In the Indian sector, positive lag correlations appear
      before February but have their maximum longitudinal extent in
      February.
    \end{enumerate}
  \end{minipage}
\end{frame}

\section{Granger Causality and Pre-Whitening Analysis}
\begin{frame}[allowframebreaks]{Granger Causality}
\tiny
% ====================================================================
% Supplemental Information: Granger Causality and Pre-Whitening
% ====================================================================
% --------------------------------------------------------------------
% SI SECTION: Methodology
% --------------------------------------------------------------------

To formally test whether Ni\~no-3.4 sea surface temperature anomalies
provide predictive information for global or tropical mean surface
temperature beyond what is contained in the temperature's own history,
we employ Granger causality testing \citep{Granger1969}.  The approach
has been previously applied in climate variability studies by, e.g.,
\citet{Mcgraw2018memory}.

Consider monthly time series $X_t$ (Ni\~no-3.4) and $Y_t$ (GMST or
tropical mean temperature).  We compare two autoregressive models for
$Y_t$.  The \emph{restricted} model uses only the history of~$Y$:
%
\begin{align}
  Y_t &= a_0 + \sum_{j=1}^{p} a_j\, Y_{t-j} + \varepsilon_t\,,
  \label{eq:granger_restricted}
\end{align}
%
while the \emph{unrestricted} model adds lagged values of~$X$:
%
\begin{align}
  Y_t &= a_0 + \sum_{j=1}^{p} a_j\, Y_{t-j}
       + \sum_{j=1}^{p} b_j\, X_{t-j} + \eta_t\,.
  \label{eq:granger_unrestricted}
\end{align}
%
In both models, $p$ is the lag order, selected by minimizing the
Bayesian Information Criterion (BIC) of the unrestricted model over
candidate values $p = 1, 2, \ldots, 18$ months.

The null hypothesis is that $X$ does not Granger-cause $Y$:
%
\begin{align}
  H_0\colon\; b_1 = b_2 = \cdots = b_p = 0\,.
  \label{eq:granger_null}
\end{align}
%
This is tested with an $F$-statistic comparing the residual sums of
squares of the restricted ($\text{RSS}_r$) and unrestricted
($\text{RSS}_u$) models:
%
\begin{align}
  F &= \frac{(\text{RSS}_r - \text{RSS}_u)\,/\,p}
            {\text{RSS}_u\,/\,(T - 2p - 1)}\,,
  \label{eq:granger_Fstat}
\end{align}
%
which follows an $F(p,\, T - 2p - 1)$ distribution under~$H_0$, where
$T$ is the effective sample size (total months minus~$p$ initial
values).  We also report the Akaike Information Criterion (AIC) and BIC
for both models as a complementary measure of whether the additional
parameters in the unrestricted model are warranted.

To test for possible bidirectional relationships or confounding, we also
perform the test in the reverse direction, asking whether the
temperature variable Granger-causes Ni\~no-3.4.

\textbf{Interpretation.}
It is important to note that Granger causality tests \emph{predictive
precedence}, not causality in the interventionist sense
\citep{Granger1969}.  The test asks whether knowledge of the past of~$X$
improves prediction of~$Y$ beyond what past~$Y$ alone provides.  In the
climate context, this is a natural and well-motivated question: we are
asking whether ENSO state carries information about the future evolution
of mean temperature that is not already encoded in the temperature's own
persistence.
\end{frame}

\begin{frame}[allowframebreaks]{Pre-Whitening and Cross-Correlation}
\tiny

As an independent check on the lag structure of the Ni\~no-3.4--temperature
relationship, we use the pre-whitening procedure of
\citet{BoxJenkins1976}.  The motivation is that if two autocorrelated
time series share a common source of low-frequency variability, their
raw cross-correlation function (CCF) will exhibit broad peaks at
multiple lags that do not correspond to any direct lagged relationship.
Pre-whitening removes this confound.

We first fit an autoregressive model of order~$q$ to~$X_t$:
%
\begin{align}
  X_t &= c + \sum_{k=1}^{q} \phi_k\, X_{t-k} + \alpha_t\,,
  \label{eq:prewhiten_ar}
\end{align}
%
where the order $q$ is selected by minimizing BIC.  The residual series
$\alpha_t = X_t - c - \sum_k \phi_k\, X_{t-k}$ is the ``pre-whitened''
version of~$X$.

The critical step is to apply the \emph{same} AR filter to~$Y_t$:
%
\begin{align}
  \beta_t &= Y_t - c - \sum_{k=1}^{q} \phi_k\, Y_{t-k}\,.
  \label{eq:prewhiten_filter_y}
\end{align}
%
Note that the coefficients $c$ and $\phi_k$ are those estimated
from~$X$, not from~$Y$.  This ensures that any cross-correlation between
$\alpha_t$ and $\beta_t$ reflects the transfer function from~$X$ to~$Y$
rather than shared autocorrelation structure \citep{BoxJenkins1976}.

The pre-whitened cross-correlation function is then
%
\begin{align}
  \hat\rho_{\alpha\beta}(\ell)
    &= \frac{1}{N}\sum_{t} \tilde\alpha_t\, \tilde\beta_{t+\ell}\,,
  \label{eq:prewhiten_ccf}
\end{align}
%
where $\tilde\alpha$ and $\tilde\beta$ are standardized to zero mean
and unit variance, and $N$ is the number of overlapping observations.
An approximate 95\% significance threshold is $\pm 2/\sqrt{N}$.

If a peak at lag~$\ell$ in the raw CCF disappears in the pre-whitened
CCF, this indicates that the raw-CCF peak was an artifact of
autocorrelation (persistence) in both series rather than evidence of a
lagged causal mechanism operating at delay~$\ell$.


\end{frame}
\begin{frame}[allowframebreaks]{Results}
\tiny
Table~\ref{tab:granger_results} summarizes the Granger causality and
pre-whitening results for both global and tropical mean temperature.
In both cases, BIC selects a lag order of $p = 2$ months.

\begin{table}[ht]\tiny
\centering
\caption{\tiny Granger causality and pre-whitening results for Ni\~no-3.4
  versus global and tropical mean surface temperature.  The lag order
  $p=2$ is selected by BIC in all cases.  $\Delta$AIC and $\Delta$BIC
  are unrestricted minus restricted (negative values favor the
  unrestricted model).  The pre-whitened CCF significance threshold is
  $\pm 0.084$.}
\label{tab:granger_results}
\begin{tabular}{l c c}
\hline\hline
  & Global & Tropical \\
\hline
\multicolumn{3}{l}{\emph{Ni\~no-3.4 $\to$ temperature}} \\
$F(2, 557)$           & 11.78               & 49.40 \\
$p$-value             & $9.7\times10^{-6}$  & $< 10^{-16}$ \\
$\Delta$AIC           & $-19.3$             & $-87.8$ \\
$\Delta$BIC           & $-10.6$             & $-79.1$ \\
$b_1$                 & $+0.069$            & $+0.110$ \\
$b_2$                 & $-0.057$            & $-0.089$ \\
\hline
\multicolumn{3}{l}{\emph{Temperature $\to$ Ni\~no-3.4}} \\
$F(2, 557)$           & 0.77                & 1.70 \\
$p$-value             & 0.46                & 0.18 \\
$\Delta$AIC           & $+2.4$              & $+0.6$ \\
$\Delta$BIC           & $+11.1$             & $+9.2$ \\
\hline
\multicolumn{3}{l}{\emph{Pre-whitened CCF (AR(2) filter)}} \\
Raw CCF at lag 3      & 0.22                & 0.53 \\
Pre-whitened CCF at lag 3 & 0.09            & $-0.01$ \\
\hline\hline
\end{tabular}
\end{table}

Ni\~no-3.4 significantly Granger-causes both global and tropical mean
temperature, with substantially stronger effects in the tropics
($F = 49.4$) than globally ($F = 11.8$).  In both cases the estimated
coefficients on lagged Ni\~no-3.4 show a short-lived impulse response:
$b_1 > 0$ (warming at lag~1) followed by $b_2 < 0$ (partial reversal
at lag~2).  This structure is consistent with a rapid atmospheric
teleconnection and is not consistent with a monotonic delayed response
peaking at lag~3.

The reverse tests find no evidence that either global ($p = 0.46$) or
tropical ($p = 0.18$) mean temperature Granger-causes Ni\~no-3.4,
consistent with the physical expectation that global- and tropical-mean
temperature does not drive tropical Pacific SST variability on monthly
timescales.

\textbf{Interpretation of $b_1 \approx -b_2$.}
When $b_1 \approx +0.07$ and $b_2 \approx -0.06$ (or $+0.11$ and
$-0.09$ for the tropics), the net effect of a step change in~$X$ that
persists for two months is approximately $b_1 + b_2 \approx 0.01$
(or $0.02$ for the tropics)---nearly zero.  The lagged $X$ terms are
acting approximately as a \emph{differencing operator} on~$X$:
%
\begin{align}
  b_1 X_{t-1} + b_2 X_{t-2} \approx b_1 (X_{t-1} - X_{t-2})
  \quad \text{when } b_2 \approx -b_1\,.
  \label{eq:granger_diff_structure}
\end{align}
%
That is, the model finds that what matters for predicting~$Y_t$ is
the \emph{change} in Ni\~no-3.4, not its level.  This maps directly
onto the tendency analysis: if Ni\~no-3.4 is steady at $+2\,^\circ$C
for several months, GMST adjusts and sits at a new equilibrium with
no further tendency.  But if Ni\~no-3.4 jumps from $+1$ to
$+2\,^\circ$C, that change forces a tendency in GMST.  Once
Ni\~no-3.4 stabilizes, the forcing on the tendency disappears.  The
$b_1 \approx -b_2$ structure is the Granger model's way of
discovering this---it is finding that the derivative of the ENSO
index, not its level, drives changes in mean temperature.

This provides a clean internal consistency across the three analyses:
the tendency correlation finds significance only at lag~0 (changes in
GMST respond to the level of ENSO contemporaneously), the Granger
model independently recovers $b_1 \approx -b_2$ (equivalent to
responding to changes in ENSO), and pre-whitening eliminates the
lag-3 peak (confirming that the level--level correlation is an
artifact of persistence).  All three are expressing the same
underlying relationship in different statistical languages.

Pre-whitening with an AR(2) filter fit to Ni\~no-3.4 ($\phi_1 = 1.47$,
$\phi_2 = -0.54$) reduces the lag-3 cross-correlation from $r = 0.22$
to $r = 0.09$ for GMST (marginally above the 95\% bound of $\pm 0.084$)
and from $r = 0.53$ to $r = -0.01$ for tropical temperature (well below
the significance bound).  The collapse of the lag-3 signal after
pre-whitening confirms that the broad peaks in the raw CCF reflect the
persistence structure of both series rather than a delayed causal
mechanism.
\end{frame}

% --------------------------------------------------------------------
% MAIN TEXT: Concise results paragraph
% --------------------------------------------------------------------

% \subsection*{Main text version}
%
% The tendency analysis (Fig.~1c,d) shows that the Ni\~no-3.4
% correlation with temperature tendencies ($\Delta Y_t = Y_t - Y_{t-1}$)
% is significant only at lag~0, for both tropical and global averages
% and regardless of detrending.  This is consistent with a
% near-instantaneous ENSO forcing on temperature tendencies rather than
% a delayed mechanism operating at the 3-month lag suggested by the raw
% cross-correlation (Fig.~1a).  A Granger causality analysis
% (supplemental information; Table~S?) confirms that Ni\~no-3.4
% provides statistically significant predictive information for both
% global and tropical mean temperature at a BIC-selected lag order of
% $p=2$ months, with substantially stronger effects in the tropics
% ($F = 49.4$, $p < 10^{-16}$) than globally ($F = 11.8$,
% $p < 10^{-5}$).  In both cases, the regression coefficients describe
% a short-lived impulse response ($b_1 > 0$, $b_2 < 0$) rather than a
% delayed response peaking at lag~3.  The reverse tests find no
% evidence that temperature Granger-causes Ni\~no-3.4.  Pre-whitening
% of both series with an AR(2) filter fit to Ni\~no-3.4 reduces the
% lag-3 cross-correlation to marginal significance for GMST and to
% effectively zero for tropical temperature, confirming that the
% apparent multi-month lag in raw cross-correlations reflects the
% persistence structure of both series.


% --------------------------------------------------------------------
% DETAILED VERSION (SI or personal notes)
% --------------------------------------------------------------------

% \subsection*{Detailed results}
%
% \paragraph{Lag order selection.}
% BIC applied to the unrestricted Granger model
% [Eq.~\eqref{eq:granger_unrestricted}] selects a lag order of $p = 2$
% months for both GMST and tropical mean temperature over the range
% $p = 1, \ldots, 18$.  For GMST, BIC at $p=2$ ($-2655.3$) is lower
% than at $p=1$ ($-2573.6$) and $p=3$ ($-2642.0$).  For tropical
% temperature, BIC at $p=2$ ($-3003.6$) is lower than at $p=1$
% ($-2962.6$) and $p=3$ ($-2998.7$).  AIC also supports $p=2$ or
% nearby values in both cases.  The consistent selection of $p=2$
% across both response variables, which is compatible with the
% $e$-folding decorrelation time of Ni\~no-3.4, indicates that the
% ENSO--temperature coupling operates on a 1--2 month timescale.
%
% \paragraph{Granger causality: Ni\~no-3.4 $\to$ temperature.}
% For GMST at $p=2$: $F(2, 557) = 11.78$, $p = 9.7 \times 10^{-6}$.
% Both AIC ($\Delta = -19.3$) and BIC ($\Delta = -10.6$) favor the
% unrestricted model.  The estimated coefficients on lagged Ni\~no-3.4
% are $b_1 = +0.069$ and $b_2 = -0.057$.
%
% For tropical mean temperature at $p=2$: $F(2, 557) = 49.40$,
% $p < 10^{-16}$.  The information criteria strongly favor the
% unrestricted model ($\Delta\text{AIC} = -87.8$;
% $\Delta\text{BIC} = -79.1$), with coefficients
% $b_1 = +0.110$ and $b_2 = -0.089$.
%
% In both cases, the positive $b_1$ and negative $b_2$ describe a
% short-lived impulse response: a positive Ni\~no-3.4 anomaly at
% time~$t$ is associated with warming at $t+1$ followed by partial
% reversal at $t+2$.  The tropical response is roughly 60\% stronger
% than the global response (comparing $b_1$ values: 0.110 vs.\ 0.069),
% consistent with dilution of the tropical ENSO signal when averaging
% over extratropical regions.  This impulse-response structure is
% consistent with the rapid atmospheric bridge mechanism
% \citep[e.g.,][]{Klein1999} and is \emph{not} consistent with a
% monotonic delayed response peaking at lag~3.
%
% \paragraph{Reverse direction: temperature $\to$ Ni\~no-3.4.}
% For GMST: $F(2, 557) = 0.77$, $p = 0.46$.  Both AIC ($\Delta = +2.4$)
% and BIC ($\Delta = +11.1$) penalize the addition of GMST lags.
% For tropical temperature: $F(2, 557) = 1.70$, $p = 0.18$.  Again
% both AIC ($\Delta = +0.6$) and BIC ($\Delta = +9.2$) penalize the
% addition.  There is no evidence that either global or tropical mean
% temperature provides predictive information for Ni\~no-3.4, consistent
% with the physical expectation that tropical Pacific SST variability is
% not driven by mean temperature on monthly timescales.  The
% unidirectional result strengthens the interpretation that the Granger
% test is detecting a genuine ENSO influence on temperature rather than
% a shared confound.
%
% \paragraph{Pre-whitening.}
% An AR(2) model fit to Ni\~no-3.4 by BIC yields coefficients
% $\phi_1 = 1.47$ and $\phi_2 = -0.54$, consistent with the damped
% oscillatory autocorrelation structure of ENSO on monthly timescales.
%
% For GMST: the raw CCF at lag~3 is $r = 0.22$.  After pre-whitening,
% this drops to $r = 0.09$, barely above the approximate 95\%
% significance threshold of $\pm 0.084$ ($= \pm 2/\sqrt{562}$).
%
% For tropical mean temperature: the raw CCF at lag~3 is $r = 0.53$,
% reflecting the much stronger ENSO--tropical temperature coupling.
% After pre-whitening, this collapses to $r = -0.01$, well below the
% significance bound.  The near-complete elimination of the lag-3
% signal in the tropics is particularly striking given the very large
% raw correlation, and demonstrates that the raw lag-3 peak was almost
% entirely attributable to shared persistence.
%
% In both cases, the lag-0 peak is essentially unchanged by
% pre-whitening, as expected since pre-whitening preserves the
% contemporaneous transfer function.
%
% \paragraph{Synthesis.}
% Three independent analyses converge on the same conclusion.
% The tendency correlation (Fig.~1c) isolates significant
% ENSO--temperature coupling only at lag~0 for both global and
% tropical averages.  The Granger test identifies genuine predictive
% power from Ni\~no-3.4 at lags 1--2 months, with an
% impulse-response structure (positive then negative coefficients)
% inconsistent with a delayed mechanism peaking at lag~3.  Pre-whitening
% reduces the lag-3 cross-correlation to marginal significance for
% GMST and to effectively zero for tropical temperature.  Together,
% these results indicate that the ENSO influence on mean temperature
% operates on a near-instantaneous to 2-month timescale, primarily
% through the tropical atmospheric bridge, and that the broad peak
% centered near lag~3 in the raw cross-correlation is predominantly an
% artifact of the strong persistence (autocorrelation) present in both
% Ni\~no-3.4 and the temperature series.  The substantially stronger
% results for tropical versus global temperature are consistent with
% this interpretation: the tropical atmosphere responds directly and
% rapidly to ENSO SST forcing, while the global signal is a diluted
% expression of this same rapid mechanism.

\begin{frame}{Accounting for seasonality: selected lag correlation maps}
  \begin{minipage}[h]{1.0\linewidth}
    \includegraphics[width=\linewidth]{/Users/tippett/notebooks/NMME/enso-t2m-figures/lag-corr-maps-tropics-month-mini.png}
  \end{minipage}
  \begin{minipage}[h]{0.6\linewidth} % Right side for the caption
    \centering \captionof{figure}{\tiny Lag correlations of
      Nino-3.4 with temperature. Each row is a different start month
      (shown on the left). The mean correlation is shown in the title
      with bold indicating the lag with maximum average value for that start
      month.}\label{}
    \begin{enumerate}\tiny
      \item Target month February is hot!
    \end{enumerate}
  \end{minipage}
\end{frame}

\begin{frame}{Accounting for seasonality: selected lag correlation maps}
  \begin{minipage}[h]{1.0\linewidth}
    \includegraphics[width=\linewidth]{/Users/tippett/notebooks/NMME/enso-t2m-figures/lag-corr-maps-tropics-month-diff-mini.png}
  \end{minipage}
  \begin{minipage}[h]{0.6\linewidth} % Right side for the caption
    \centering \captionof{figure}{\tiny Lag correlations of Nino-3.4
      with temperature \textbf{tendencies}. Each row is a different
      start month (shown on the left). The mean correlation is shown
      in the title with bold indicating the lag with maximum average
      value for that start month.}\label{}
    \begin{enumerate}\tiny
    \item Lag 1 or zero is where the warming happens.
      \item So we look for the mechanisms at lag zero.
    \end{enumerate}
  \end{minipage}
\end{frame}

\begin{frame}{Accounting for seasonality: Atlantic, Indian, and south China Sea sectors}
  \begin{minipage}[h]{0.38\linewidth}
    \includegraphics[width=\linewidth]{/Users/tippett/notebooks/NMME/enso-t2m-figures/lag-corr-reg-with-seasonality-2-lines-atl-ind.pdf}
  \end{minipage}\hspace{0.5em}
  \begin{minipage}[h]{0.5\linewidth} % Right side for the caption
    \centering \captionof{figure}{\tiny Lag correlations and
      regression coefficients as a function of start month and target
      month of Nino-3.4 with (a, c) tropical Atlantic sector mean temperature
      and (b, d) Indian sector mean temperature.  Each line correspond
      to a different start month. Gray (black) dots indicate
      statistical significance at the 5\% (1\%)
      level.}\label{fig:corr-reg-with-seasonality-regions}
    \begin{enumerate}\tiny
    \item Both peak in the target month of February but the Atlantic
      peak is narrow.
    \end{enumerate}
  \end{minipage}
\end{frame}


\begin{frame}{Accounting for seasonality}
  \begin{minipage}[t]{0.5\linewidth}
    \vspace{-15pt}
    \includegraphics[width=\linewidth]{/Users/tippett/notebooks/NMME/enso-t2m-figures/lag-corr-zonal-mean-june.pdf}
    \includegraphics[width=\linewidth]{/Users/tippett/notebooks/NMME/enso-t2m-figures/lag-corr-meridional-mean-tropics-june.pdf}
    \centering \captionof{figure}{\fontsize{5pt}{7pt}\selectfont
      Month-dependent lag correlations of Nino-3.4 with zonal mean
      temperature. The start month is shown in the title of each
      panel, and the horizontal axis corresponds to the target
      month. Colored vertical lines in selected panels (see text) mark
      Oct/Nov (purple), Jan/Feb (green), and May/Jun (blue)
      boundaries. Dashed lines are at 22.5\textdegree N and
      22.5\textdegree S. White (black) dots indicate statistical
      significance at the 5\% (1\%) level.}\label{fig:june-starts-corr}
  \end{minipage}\hfill
  \begin{minipage}[t]{0.49\linewidth} % Right side for the caption
    \vspace{-10pt}
    \begin{enumerate}\fontsize{5pt}{7pt}\selectfont
    \item To see the meridional dependence of the relation between
      ENSO and surface temperature, we computed the month-dependent
      lag correlation of Nino-3.4 with zonal mean temperature for all
      start months and found they were all similar to June starts
      (Fig.\ \ref{fig:june-starts-corr}a).
    \item Correlations are mostly confined to the tropics except for
      some weak correlations in the southern hemisphere. Ref?
    \item The dominant feature is tropics-wide lag correlations that
      weaken and fade when the target month reaches May/June (blue
      lines) for start months June--March. This is the ENSO spring
      persistence barrier.
    \item Within the tropics, another seasonal feature is positive
      correlations that expand southward in November (purple lines)
      and persist through the following summer for May--October start
      months.
    \item For June--December start months, positive lag
      correlations in the tropics expand northward in February (green
      lines), though with little persistence.
    \item Zonal dependence. Correlations in the tropical Pacific hit
      the spring persistence barrier around June.  For May--October
      start months, positive lag correlations in the tropical Indian
      sector light up in November and persist through the following
      summer.  Further east (south China Sea?), the action starts
      later.
    \item Seasonal feature 3: For June--December start months,
      positive lag correlations in the tropical Atlantic expand
      northward in February and don't really persist.
    \item For June--December start months, positive lag correlations
      appear in the Atlantic sector beginning in February.
    \item In the Indian sector, positive lag correlations appear
      before February but have their maximum longitudinal extent in
      February.
    \end{enumerate}
  \end{minipage}
\end{frame}

\begin{frame}{Accounting for seasonality: Zonal means}
  \begin{minipage}[t]{0.55\linewidth}
    \vspace{-15pt}
    \includegraphics[width=\linewidth]{/Users/tippett/notebooks/NMME/enso-t2m-figures/lag-corr-zonal-mean.pdf}
    \centering \captionof{figure}{\tiny Month-dependent lag correlations of
      Nino-3.4 with zonal mean temperature. The start month is shown
      in the title of each panel, and the horizontal axis corresponds
      to the target month. Colored vertical lines in selected panels
      (see text) mark Oct/Nov (purple), Jan/Feb (green), and May/Jun
      (blue) boundaries. Dashed lines are at 22.5\textdegree N and
      22.5\textdegree S. White (black) dots indicate statistical
      significance at the 5\% (1\%) level. }\label{fig:zonal-means}
  \end{minipage}\hfill
  \begin{minipage}[t]{0.45\linewidth} % Right side for the caption
    \vspace{0pt}
    \begin{enumerate}\tiny
    \item To see the meridional dependence of the  relation
      between ENSO and surface temperature, we computed the
      month-dependent lag correlation of Nino-3.4 with zonal mean
      temperature (Fig.\ \ref{fig:zonal-means}).
    \item Correlations are mostly confined to the tropics except for
      some weak correlations in the southern hemisphere. Ref?
    \item The dominant feature is tropics-wide lag correlations that
      weaken and fade when the target month reaches May/June (blue
      lines) for start months June--March. This is the ENSO spring
      persistence barrier.
    \item Within the tropics, another seasonal feature is positive
      correlations that expand southward in November (purple lines)
      and persist through the following summer for May--October start
      months.
    \item For June--December start months, positive lag
      correlations in the tropics expand northward in February (green
      lines), though with little persistence.
    \end{enumerate}
  \end{minipage}
\end{frame}


\begin{frame}{Accounting for seasonality: Meridional means}
  \begin{minipage}[t]{0.55\linewidth}
    \vspace{-15pt}
    \includegraphics[width=\linewidth]{/Users/tippett/notebooks/NMME/enso-t2m-figures/lag-corr-meridional-mean-tropics.pdf}
    \centering \captionof{figure}{\tiny Lag correlations and
      regressions of Nino-3.4 with tropical meridional average
      temperature. White (black) dots indicate statistical
      significance at the 5\% (1\%) level.}\label{fig:meridional-means-cor}
  \end{minipage}\hfill
  \begin{minipage}[t]{0.45\linewidth} % Right side for the caption
    \vspace{0pt}
    \begin{enumerate}\tiny
    \item Seasonal feature 1: Correlations in the tropical Pacific hit
      the spring persistence barrier around June.

    \item Seasonal feature 2: For May--October start months, positive
      lag correlations in the tropical Indian sector light up in
      November and persist through the following summer.
    \item Further east (south China Sea?), the action starts later.
    \item Seasonal feature 3: For June--December start months,
      positive lag correlations in the tropical Atlantic expand
      northward in February and don't really persist.
    \item For June--December start months, positive lag
      correlations appear in the Atlantic sector beginning in
      February.
    \item In the Indian sector, positive lag correlations appear
      before February but have their maximum longitudinal extent in
      February.
    \end{enumerate}
  \end{minipage}
\end{frame}

\begin{frame}{Accounting for seasonality: Meridional means}
  \begin{minipage}[t]{0.55\linewidth}
    \vspace{-15pt}
    \includegraphics[width=\linewidth]{/Users/tippett/notebooks/NMME/enso-t2m-figures/lag-reg-meridional-mean-tropics.pdf}
    \centering \captionof{figure}{\tiny Lag correlations and
      regressions of Nino-3.4 with tropical meridional average
      temperature. White (black) dots indicate statistical
      significance at the 5\% (1\%) level.}\label{fig:meridional-means-reg}
  \end{minipage}\hfill
  \begin{minipage}[t]{0.45\linewidth} % Right side for the caption
    \vspace{0pt}
    \begin{enumerate}\tiny
    \item Seasonal feature 1: Correlations in the tropical Pacific hit
      the spring persistence barrier around June.

    \item Seasonal feature 2: For May--October start months, positive
      lag correlations in the tropical Indian sector light up in
      November and persist through the following summer.
    \item Further east (south China Sea?), the action starts later.
    \item Seasonal feature 3: For June--December start months,
      positive lag correlations in the tropical Atlantic expand
      northward in February and don't really persist.
    \item For June--December start months, positive lag
      correlations appear in the Atlantic sector beginning in
      February.
    \item In the Indian sector, positive lag correlations appear
      before February but have their maximum longitudinal extent in
      February.
    \end{enumerate}
  \end{minipage}
\end{frame}

\section{Leftovers}
\section{Bridges}

\begin{frame}{Atlantic: Cloud–shortwave feedback}
  \begin{minipage}[h]{\linewidth}
    \includegraphics[width=0.49\linewidth]{/Users/tippett/notebooks/NMME/enso-t2m-figures/map-total-cloud-cover-n34-corr-atl.png}
    \includegraphics[width=0.49\linewidth]{/Users/tippett/notebooks/NMME/enso-t2m-figures/map-surface-net-solar-radiation-n34-corr-atl.png}
  \end{minipage}\hspace{0.5em}
  
  \begin{minipage}[h]{0.8\linewidth} % Right side for the caption
    \centering \captionof{figure}{\tiny }
    \begin{enumerate}\tiny
    \item 
    \end{enumerate}
  \end{minipage}
\end{frame}

\begin{frame}{Atlantic: Walker circulation anomalies}
  \begin{minipage}[h]{\linewidth}
    \includegraphics[width=0.49\linewidth]{/Users/tippett/notebooks/NMME/enso-t2m-figures/map-U850-n34-corr-atl.png}
    \includegraphics[width=0.49\linewidth]{/Users/tippett/notebooks/NMME/enso-t2m-figures/map-850-windspeed-n34-corr-atl.png}
  \end{minipage}\hspace{0.5em}
  
  \begin{minipage}[h]{0.8\linewidth} % Right side for the caption
    \centering \captionof{figure}{\tiny }
    \begin{enumerate}\tiny
    \item 
    \end{enumerate}
  \end{minipage}
\end{frame}

\begin{frame}{Atlantic: Hadley cell modulation}
  \begin{minipage}[h]{\linewidth}
    \includegraphics[width=0.49\linewidth]{/Users/tippett/notebooks/NMME/enso-t2m-figures/map-V850-n34-corr-atl.png}
    \includegraphics[width=0.49\linewidth]{/Users/tippett/notebooks/NMME/enso-t2m-figures/map-MSLP-n34-corr-atl.png}
  \end{minipage}\hspace{0.5em}
  
  \begin{minipage}[h]{0.8\linewidth} % Right side for the caption
    \centering \captionof{figure}{\tiny }
    \begin{enumerate}\tiny
    \item 
    \end{enumerate}
  \end{minipage}
\end{frame}

\begin{frame}{Atlantic: Wind–Evaporation–SST feedback}
  \begin{minipage}[h]{\linewidth}
    \includegraphics[width=0.32\linewidth]{/Users/tippett/notebooks/NMME/enso-t2m-figures/map-U850-n34-corr-atl.png}
    \includegraphics[width=0.32\linewidth]{/Users/tippett/notebooks/NMME/enso-t2m-figures/map-V850-n34-corr-atl.png}
    \includegraphics[width=0.32\linewidth]{/Users/tippett/notebooks/NMME/enso-t2m-figures/map-surface-latent-heat-flux-n34-corr-atl.png}    
  \end{minipage}\hspace{0.5em}
  
  \begin{minipage}[h]{0.8\linewidth} % Right side for the caption
    \centering \captionof{figure}{\tiny }
    \begin{enumerate}\tiny
    \item 
    \end{enumerate}
  \end{minipage}
\end{frame}


\begin{frame}{Atlantic 2}
  \begin{minipage}[h]{\linewidth}
    \includegraphics[width=0.49\linewidth]{/Users/tippett/notebooks/NMME/enso-t2m-figures/map-850winds-n34-corr-atl.png}
    \includegraphics[width=0.49\linewidth]{/Users/tippett/notebooks/NMME/enso-t2m-figures/map-surface-latent-heat-flux-n34-corr-atl.png}
  \end{minipage}\hspace{0.5em}
  
  \begin{minipage}[h]{0.8\linewidth} % Right side for the caption
    \centering \captionof{figure}{\tiny Correlation of surface net solar radiation with
      Nino-3.4. Only correlations greater than}
    \begin{enumerate}\tiny
    \item Winds go with evaporation.
    \end{enumerate}
  \end{minipage}
\end{frame}

\begin{frame}{Atlantic 2}
  \begin{minipage}[h]{\linewidth}
    \includegraphics[width=0.32\linewidth]{/Users/tippett/notebooks/NMME/enso-t2m-figures/map-omega500-n34-corr-atl.png}
    \includegraphics[width=0.32\linewidth]{/Users/tippett/notebooks/NMME/enso-t2m-figures/map-surface-net-solar-radiation-n34-corr-atl.png}
    \includegraphics[width=0.32\linewidth]{/Users/tippett/notebooks/NMME/enso-t2m-figures/map-total-cloud-cover-n34-corr-atl.png}
  \end{minipage}\hspace{0.5em}
  
  \begin{minipage}[h]{0.8\linewidth} % Right side for the caption
    \centering \captionof{figure}{\tiny Correlation of surface net solar radiation with
      Nino-3.4. Only correlations greater than}
    \begin{enumerate}\tiny
    \item Subsidence (positive vertical velocity) goes with more net
      solar radiation and cloudiness.
    \end{enumerate}
  \end{minipage}
\end{frame}


\begin{frame}{Indo-western Pacific: Cloud–shortwave feedback}
  \begin{minipage}[h]{\linewidth}
    \includegraphics[width=0.49\linewidth]{/Users/tippett/notebooks/NMME/enso-t2m-figures/map-total-cloud-cover-n34-corr-indo.png}
    \includegraphics[width=0.49\linewidth]{/Users/tippett/notebooks/NMME/enso-t2m-figures/map-surface-net-solar-radiation-n34-corr-indo.png}
  \end{minipage}\hspace{0.5em}
  
  \begin{minipage}[h]{0.8\linewidth} % Right side for the caption
    \centering \captionof{figure}{\tiny }
    \begin{enumerate}\tiny
    \item Reduced convective cloud during El Ni\~no increases
      insolation;
    \end{enumerate}
  \end{minipage}
\end{frame}

\begin{frame}{Indo-western Pacific: Walker circulation anomalies}
  \begin{minipage}[h]{\linewidth}
    \includegraphics[width=0.49\linewidth]{/Users/tippett/notebooks/NMME/enso-t2m-figures/map-U850-n34-corr-indo.png}
    \includegraphics[width=0.49\linewidth]{/Users/tippett/notebooks/NMME/enso-t2m-figures/map-850-windspeed-n34-corr-indo.png}
  \end{minipage}\hspace{0.5em}
  
  \begin{minipage}[h]{0.8\linewidth} % Right side for the caption
    \centering \captionof{figure}{\tiny }
    \begin{enumerate}\tiny
    \item Weakened ascending/descending branches during El Ni\~no
    \end{enumerate}
  \end{minipage}
\end{frame}

\begin{frame}{Indo-western Pacific: Hadley cell modulation}
  \begin{minipage}[h]{\linewidth}
    \includegraphics[width=0.49\linewidth]{/Users/tippett/notebooks/NMME/enso-t2m-figures/map-V850-n34-corr-indo.png}
    \includegraphics[width=0.49\linewidth]{/Users/tippett/notebooks/NMME/enso-t2m-figures/map-MSLP-n34-corr-indo.png}
  \end{minipage}\hspace{0.5em}
  
  \begin{minipage}[h]{0.8\linewidth} % Right side for the caption
    \centering \captionof{figure}{\tiny }
    \begin{enumerate}\tiny
    \item 
    \end{enumerate}
  \end{minipage}
\end{frame}

\begin{frame}{Indo-western Pacific: Wind–Evaporation–SST feedback}
  \begin{minipage}[h]{\linewidth}
    \includegraphics[width=0.32\linewidth]{/Users/tippett/notebooks/NMME/enso-t2m-figures/map-U850-n34-corr-indo.png}
    \includegraphics[width=0.32\linewidth]{/Users/tippett/notebooks/NMME/enso-t2m-figures/map-V850-n34-corr-indo.png}
    \includegraphics[width=0.32\linewidth]{/Users/tippett/notebooks/NMME/enso-t2m-figures/map-surface-latent-heat-flux-n34-corr-indo.png}    
  \end{minipage}\hspace{0.5em}
  
  \begin{minipage}[h]{0.8\linewidth} % Right side for the caption
    \centering \captionof{figure}{\tiny }
    \begin{enumerate}\tiny
    \item 
    \end{enumerate}
  \end{minipage}
\end{frame}

\begin{frame}{Indo-western Pacific 2}
  \begin{minipage}[h]{\linewidth}
    \includegraphics[width=0.49\linewidth]{/Users/tippett/notebooks/NMME/enso-t2m-figures/map-850winds-n34-corr-indo.png}
    \includegraphics[width=0.49\linewidth]{/Users/tippett/notebooks/NMME/enso-t2m-figures/map-surface-latent-heat-flux-n34-corr-indo.png}
  \end{minipage}\hspace{0.5em}
  
  \begin{minipage}[h]{0.8\linewidth} % Right side for the caption
    \centering \captionof{figure}{\tiny Correlation of surface net solar radiation with
      Nino-3.4. Only correlations greater than}
    \begin{enumerate}\tiny
    \item Winds go with evaporation.
    \end{enumerate}
  \end{minipage}
\end{frame}

\begin{frame}{Indo-western Pacific 2}
  \begin{minipage}[h]{\linewidth}
    \includegraphics[width=0.32\linewidth]{/Users/tippett/notebooks/NMME/enso-t2m-figures/map-omega500-n34-corr-indo.png}
    \includegraphics[width=0.32\linewidth]{/Users/tippett/notebooks/NMME/enso-t2m-figures/map-surface-net-solar-radiation-n34-corr-indo.png}
    \includegraphics[width=0.32\linewidth]{/Users/tippett/notebooks/NMME/enso-t2m-figures/map-total-cloud-cover-n34-corr-indo.png}
  \end{minipage}\hspace{0.5em}
  
  \begin{minipage}[h]{0.8\linewidth} % Right side for the caption
    \centering \captionof{figure}{\tiny Correlation of surface net solar radiation with
      Nino-3.4. Only correlations greater than}
    \begin{enumerate}\tiny
    \item Subsidence (positive vertical velocity) goes with more net
      solar radiation and cloudiness.
    \end{enumerate}
  \end{minipage}
\end{frame}


\begin{frame}{Explaining the seasonality}
  The presence of seasonality in the month-dependent lag correlations
  raises the questions regarding its origin. Namely,
  \begin{enumerate}
  \item Where in latitude is the source of seasonality? A: MostlyThe
    tropics. (Fig.\ \ref{fig:zonal-means})
  \item Q: When in the tropics? A: Things happen in November,
    February, and May (Fig.\ \ref{fig:zonal-means}).
  \item Q: What happens in November, February, and May? A: Indian
    sector in November, Atlantic in February, and Pacific in
    May. (Fig.\ \ref{fig:meridional-means})
  \end{enumerate}
\end{frame}

\begin{frame}{Making sense of the seasonality}
  \begin{enumerate}
  \item Averaging the month-dependent lag correlation matches the
    all-months lag correlation (Fig.\ \ref{fig:pooled-vs-ave}).
  \item One explanation: composites for warm and cool ENSO show
    seasonality in the tropics (Figs\ \ref{fig:composite}). Composites
    are easy to understand but don't directly show the seasonality in
    correlation.
  \item On the other hand, directly show the seasonality in
    correlation: month-dependent lag correlations calculated as in
    Fig.\ \ref{fig:with-seasonality} but for zonal means, tropical
    mean, and maps (Figs.\ \ref{fig:june-starts-corr} and
    \ref{fig:june-starts-reg}). Or boxes (Fig.\
    \ref{fig:corr-reg-with-seasonality-regions}).
  \item Budget (regression coefficients for regions weighted by their
    size)?
  \item Bridge/tunnel check (Figs.\ 11--14)?
  \item Implications?
  \item The right model? $T_\tau = T_{tau-1} + N_\tau + \dots $
  \item Idea: does using the right model reduce trend uncertainty? No.
  \end{enumerate}
\end{frame}


\begin{frame}{No accounting for seasonality or taste}
  \begin{minipage}[h]{0.65\linewidth}
    \includegraphics[width=\linewidth]{/Users/tippett/notebooks/NMME/enso-t2m-figures/lag-corr-reg-without-seasonality-mini.pdf}
  \end{minipage}\hspace{0.25em}
  \begin{minipage}[h]{0.32\linewidth} % Right side for the caption
    \centering \captionof{figure}{\tiny Lag (a) correlations and
      (b) regression coefficients of Nino-3.4 with  global   and
      tropical mean temperature, (b, d) global and tropical mean
      temperature tendencies.  Filled circles indicate statistical
      significance at the 1\% level.  Dark colors indicate that both
      time series are detrended.}\label{fig:no-seasonality-corr-reg-diff}
    \begin{enumerate}\tiny
    \item (a, c) Correlation and regression coefficients peak at lags
      of about three or four months for global and tropical mean
      temperatures. Correlations of Nino-3.4 with global mean
      temperature are substantially weaker than those with tropical
      mean temperature.
    \item Detrending increases correlations. The impact of detrending
      on regression coefficients is more modest since trends in
      Nino-3.4 are modest.
    \item Detrending does nothing to tendency correlations and
      regressions.
    \item Correlations with tendencies peak at lag zero.  The only
      statistically significant correlations with tendencies are those
      with the tropical mean at lags zero and one (warming) and at lag
      six and greater (cooling).
    \end{enumerate}
  \end{minipage}
\end{frame}


\begin{frame}{Accounting for seasonality: Zonal mean tendencies}
  \begin{minipage}[h]{0.7\linewidth}
    \includegraphics[width=\linewidth]{/Users/tippett/notebooks/NMME/enso-t2m-figures/lag-corr-zonal-mean-diff.pdf}
  \end{minipage}
  \begin{minipage}[h]{0.29\linewidth} % Right side for the caption
    \centering \captionof{figure}{\tiny Lag correlations as a
      function of start month and lag of Nino-3.4 with temperature
      \textbf{tendencies}. White (black) dots indicate statistical
      significance at the 5\% (1\%)
      level. }\label{fig:zonal-means-diff}
    \begin{enumerate}\tiny
      \item Cooling on the equator sets in during February for starts months June--January.
      \item Tropical expansion southward of warming in November for
        start months May--October.
    \item Tropical expansion northward of weak warming in February for
      start months June--December.
    \end{enumerate}
  \end{minipage}
\end{frame}

\begin{frame}{Accounting for seasonality: Meridional mean tendencies}
  \begin{minipage}[h]{0.7\linewidth}
    \includegraphics[width=\linewidth]{/Users/tippett/notebooks/NMME/enso-t2m-figures/lag-corr-meridional-mean-tropics-diff.pdf}
  \end{minipage}
  \begin{minipage}[h]{0.29\linewidth} % Right side for the caption
    \centering \captionof{figure}{\tiny Lag correlations and
      regressions of Nino-3.4 with tropical meridional average
      temperature. White (black) dots indicate statistical
      significance at the 5\% (1\%) level.}\label{fig:meridional-means}
    \begin{enumerate}\tiny
    \item ???
    \end{enumerate}
  \end{minipage}
\end{frame}


\begin{frame}{Nino-3.4 correlation with U850}
  \begin{minipage}[h]{0.75\linewidth}
    \includegraphics[width=\linewidth]{/Users/tippett/notebooks/NMME/enso-t2m-figures/map-U850-n34-corr.png}
  \end{minipage}
  \begin{minipage}[h]{0.24\linewidth} % Right side for the caption
    \centering \captionof{figure}{\tiny Correlation of U850 with Nino-3.4. Only correlations greater than}
    \begin{enumerate}\tiny
    \item 
    \end{enumerate}
  \end{minipage}
\end{frame}

\begin{frame}{Nino-3.4 correlation with V850}
  \begin{minipage}[h]{0.75\linewidth}
    \includegraphics[width=\linewidth]{/Users/tippett/notebooks/NMME/enso-t2m-figures/map-V850-n34-corr.png}
  \end{minipage}
  \begin{minipage}[h]{0.24\linewidth} % Right side for the caption
    \centering \captionof{figure}{\tiny Correlation of V850 with Nino-3.4. Only correlations greater than}
    \begin{enumerate}\tiny
    \item 
    \end{enumerate}
  \end{minipage}
\end{frame}

\begin{frame}{Nino-3.4 correlation with 850 windspeed}
  \begin{minipage}[h]{0.75\linewidth}
    \includegraphics[width=\linewidth]{/Users/tippett/notebooks/NMME/enso-t2m-figures/map-850-windspeed-n34-corr.png}
  \end{minipage}  
  \begin{minipage}[h]{0.24\linewidth} % Right side for the caption
    \centering \captionof{figure}{\tiny Correlation of 850 windspeed with Nino-3.4. Only correlations greater than}
    \begin{enumerate}\tiny
    \item 
    \end{enumerate}
  \end{minipage}
\end{frame}

\begin{frame}{Nino-3.4 correlation with surface latent heat flux}
  \begin{minipage}[h]{0.75\linewidth}
    \includegraphics[width=\linewidth]{/Users/tippett/notebooks/NMME/enso-t2m-figures/map-surface-latent-heat-flux-n34-corr.png}
  \end{minipage}  
  \begin{minipage}[h]{0.24\linewidth} % Right side for the caption
    \centering \captionof{figure}{\tiny Winds go with evaporation.}
    \begin{enumerate}\tiny
    \item 
    \end{enumerate}
  \end{minipage}
\end{frame}

\begin{frame}{Nino-3.4 correlation with omega 500}
  \begin{minipage}[h]{0.75\linewidth}
    \includegraphics[width=\linewidth]{/Users/tippett/notebooks/NMME/enso-t2m-figures/map-omega500-n34-corr.png}
  \end{minipage}  
  \begin{minipage}[h]{0.24\linewidth} % Right side for the caption
    \centering \captionof{figure}{\tiny Subsidence goes with more net
      solar radiation. (Positive omega 500 values are downward).}
    \begin{enumerate}\tiny
    \item 
    \end{enumerate}
  \end{minipage}
\end{frame}

\begin{frame}{Nino-3.4 correlation with surface net solar radiation}
  \begin{minipage}[h]{0.75\linewidth}
    \includegraphics[width=\linewidth]{/Users/tippett/notebooks/NMME/enso-t2m-figures/map-surface-net-solar-radiation-n34-corr.png}
  \end{minipage}  
  \begin{minipage}[h]{0.24\linewidth} % Right side for the caption
    \centering \captionof{figure}{\tiny Correlation of surface net solar radiation with
      Nino-3.4. Only correlations greater than}
    \begin{enumerate}\tiny
    \item 
    \end{enumerate}
  \end{minipage}
\end{frame}


\begin{frame}{Grainger causality analysis}
  \begin{enumerate}\footnotesize
  \item Granger causality analysis offers advantages over lag
    regression by taking into account the autocorrelation of the data
    \citep{Mcgraw2018memory}.
  \item More precisely, Grainger causality analysis addresses the
    question of whether predictions of $Y$ based on past values of $X$
    and $Y$ are better than predictions of $Y$ based on past values of
    $Y$ alone. Thus, strictly speaking, predictive ability is the
    issue addressed rather than causality.
  \item Here $X$ is the Nino-3.4 index and $Y$ is either tropical mean
    temperature or global mean temperature.
  \item Typically for non-stationary variables (such as average temperatures
    with their upward trends), the data are differenced (until
    stationary), and the analysis proceeds with the differenced data.
  \item We differenced the global and tropical mean temperature data and checked
    that no auto correlation was left (Fig.\ \ref{fig:data-acfs}),
    which supported using no past  values of the differenced data in
    the analysis.
  \item We then tested the statistical significance of the
    regression coefficient $b_p$ in
    \begin{equation*}
      Y_t - Y_{t-1} = b_p X_{t-p}\,, \qquad p=0, 1, \dots
    \end{equation*}
  \item In practical terms, we tested whether predictions of changes
    in global and tropical mean temperature are improved by knowing
    past values of Nino-3.4.
  \end{enumerate}
  
\end{frame}

\section{No accounting for seasonality}

\begin{frame}%{No accounting for seasonality}
  \begin{minipage}[h]{0.75\linewidth}
    \includegraphics[width=\linewidth]{/Users/tippett/notebooks/NMME/enso-t2m-figures/lag-reg-maps.png}
  \end{minipage}
  \begin{minipage}[h]{0.24\linewidth} % Right side for the caption
    \centering \captionof{figure}{Lag \textbf{regression} of Nino-3.4 with
      2m temperature. White (black) dots indicate statistical
      significance at the 5\% (1\%) level. The global mean of the
      correlation is shown in the title with bold indicating the lag
      with maximum value.}\label{fig:no-seasonality-zonal}
  \end{minipage}
\end{frame}

\begin{frame}
  \begin{enumerate}
  \item The maximum global averaged correlation occurs at lag three.
  \item At lag three, the correlation in the tropical Pacific has
    decreased, which \citet[][see their Fig. 2]{Klein1999remote} said is an indication
    that correlations in other basins reach their peak as the
    ENSO-related correlations in the Pacific are decaying.
  \end{enumerate}
\end{frame}

\begin{frame}%{No accounting for seasonality}
  \begin{minipage}[h]{0.75\linewidth}
    \includegraphics[width=\linewidth]{/Users/tippett/notebooks/NMME/enso-t2m-figures/lag-reg-maps-diff.png}
  \end{minipage}
  \begin{minipage}[h]{0.24\linewidth} % Right side for the caption
    \centering \captionof{figure}{Lag \textbf{regression} of Nino-3.4 with
      temperature \textbf{tendencies}. White (black) dots indicate
      statistical significance at the 5\% (1\%) level. The global mean
      of the correlation is shown in the title with bold indicating
      the lag with maximum value.}\label{fig:no-seasonality-zonal}
  \end{minipage}
\end{frame}

\begin{frame}
  \begin{enumerate}
  \item The maximum global averaged correlation with
    \textbf{tendencies} occurs at lag zero. Mostly over ocean except
    for the Amazon (the ``Green Ocean'' especially during Dec--May wet
    season).
  \item As the lag increases, correlations in the tropical Pacific
    become negative (ENSO decay) while remaining positive in tropical
    north Atlantic (closer to South America) through lag 3, Northwest
    Indian ocean through lag 2, and western tropical Pacific through
    lag 7 or so.
  \end{enumerate}
\end{frame}


\begin{frame}%{No accounting for seasonality}
  \begin{minipage}[h]{0.75\linewidth}
    \includegraphics[width=\linewidth]{/Users/tippett/notebooks/NMME/enso-t2m-figures/zonal-lag-corr-reg-without-seasonality.pdf}
  \end{minipage}
  \begin{minipage}[h]{0.24\linewidth} % Right side for the caption
    \centering \captionof{figure}{Lag correlations and regressions
      of Nino-3.4 with zonal mean temperature. White (black) dots
      indicate statistical significance at the 5\% (1\%) level. Dashed
      lines are at $\pm$
      27.5\textdegree.}\label{fig:no-seasonality-zonal}
  \end{minipage}
\end{frame}

\begin{frame}%{No accounting for seasonality}
  \begin{minipage}[h]{0.75\linewidth}
    \includegraphics[width=\linewidth]{/Users/tippett/notebooks/NMME/enso-t2m-figures/zonal-lag-corr-reg-without-seasonality-diff.pdf}
  \end{minipage}\hspace{0.25em}
  \begin{minipage}[h]{0.23\linewidth} % Right side for the caption
    \centering \captionof{figure}{Lag correlations and regressions
      of Nino-3.4 with zonal mean temperature
      \textbf{tendencies}. White (black) dots indicate statistical
      significance at the 5\% (1\%) level. Dashed lines are at $\pm$
      27.5\textdegree.}\label{fig:no-seasonality-zonal}
  \end{minipage}
\end{frame}

\begin{frame}
  \begin{enumerate}
  \item More ways of seeing that nothing gets outside the tropics (on average).
  \item More ways of seeing that there is no need to detrend when
    looking at tendencies.
  \end{enumerate}
\end{frame}

\begin{frame}%{No accounting for seasonality}
  \begin{minipage}[h]{0.65\linewidth}
    \includegraphics[width=\linewidth]{/Users/tippett/notebooks/NMME/enso-t2m-figures/lag-corr-reg-without-seasonality.pdf}
  \end{minipage}\hspace{0.25em}
  \begin{minipage}[h]{0.32\linewidth} % Right side for the caption
    \centering \captionof{figure}{Lag correlations of Nino-3.4 with
      (a) global, tropical, (b) extratropical, NH extratropical, and
      SH extratropical average temperature. Lag regressions between
      Nino-3.4 and (c) global, tropical, (d) extratropical, NH
      extratropical, and SH extratropical average temperature. Filled
      circles indicate statistical significance at the 5\% level.
      Light colors indicate that both time series are
      detrended.}\label{fig:no-seasonality-plus}
  \end{minipage}
\end{frame}

\begin{frame}
  \begin{enumerate}
  \item The correlation with tropical and global mean surface temperature is
    highest at lags three and four (Fig.\ \ref{fig:no-seasonality}a).
  \item Detrending increases correlations.
  \item Detrending results in little change in the regression
    coefficients for tropical and global mean surface temperature
    (Fig.\ \ref{fig:no-seasonality}c) because the trend in Nino-3.4 is
    small and only its variance appears in the regression coefficient
    (in the denominator).
  \item Relations with extratropical mean temperatures are much weaker
    (Fig.\ \ref{fig:no-seasonality}b, d). Only a two correlations with
    NH extratropical temperature are statistically
    significant. Multiple testing.
  \item By pooling all months together, all this analysis describes
    ``on average'' behavior with the potential for different behavior
    in different years and in different months of the year.
  \end{enumerate}
\end{frame}

\begin{frame}%{No accounting for seasonality}
  \begin{minipage}[h]{0.65\linewidth}
    \includegraphics[width=\linewidth]{/Users/tippett/notebooks/NMME/enso-t2m-figures/lag-corr-reg-without-seasonality-diff.pdf}
  \end{minipage}\hspace{0.5em}
  \begin{minipage}[h]{0.32\linewidth} % Right side for the caption
    \centering \captionof{figure}{Lag correlations of Nino-3.4 with
      \textbf{tendencies} of (a) global, tropical, (b) extratropical,
      NH extratropical, and SH extratropical average
      temperature. Lag regressions between Nino-3.4 and (c) global,
      tropical, (d) extratropical, NH extratropical, and SH
      extratropical average temperature. Filled circles indicate
      statistical significance at the 5\% level.  Light colors
      indicate that both time series are
      detrended.}\label{fig:no-seasonality-diff}
  \end{minipage}
\end{frame}

\begin{frame}
  \begin{enumerate}
  \item The correlation with tropical and global mean surface
    \textbf{tendencies} temperature is highest at zero lag (Fig.\
    \ref{fig:no-seasonality-diff}a). (There is some problem with the
    significance).
  \item Detrending does nothing because differencing removes the
    trend?
  \item Relations with extratropical mean temperatures are much weaker
    (Fig.\ \ref{fig:no-seasonality-diff}b, d).
  \end{enumerate}
\end{frame}

\begin{frame}%{No accounting for seasonality}
  \begin{minipage}[h]{0.75\linewidth}
    \includegraphics[width=\linewidth]{/Users/tippett/notebooks/NMME/enso-t2m-figures/lag-corr-meridional-mean-tropics-with-map.pdf}
  \end{minipage}
  \begin{minipage}[h]{0.24\linewidth} % Right side for the caption
    \centering \captionof{figure}{Lag correlations and regressions
      of Nino-3.4 with the mean of -27.5 to 27.5 mean temperature and
      temperature \textbf{tendencies}. White (black) dots indicate
      statistical significance at the 5\% (1\%) level. With reference
      maps for the geographically challenged.}\label{}
  \end{minipage}
\end{frame}

\section{Accounting for seasonality}
\begin{frame}{Global temperature}
  \begin{minipage}[h]{0.6\linewidth}
    \includegraphics[width=\linewidth]{/Users/tippett/notebooks/NMME/enso-t2m-figures/lag-corr-maps-month.png}
  \end{minipage}\hfill
  \begin{minipage}[h]{0.28\linewidth} % Right side for the caption
    \centering \captionof{figure}{Lag correlations as a function of
      start month and lag of Nino-3.4 with temperature. The global
      mean of the correlation is shown in the title with bold
      indicating the lag with maximum value.
      % White (black) dots indicate statistical significance at the
      % 5\% (1\%) level.
    }\label{fig:corr-with-seasonality-2}
    \begin{enumerate}\scriptsize
      \item 
    \end{enumerate}
  \end{minipage}
\end{frame}

\begin{frame}{Global temperature tendencies}
  \begin{minipage}[h]{0.6\linewidth}
    \includegraphics[width=\linewidth]{/Users/tippett/notebooks/NMME/enso-t2m-figures/lag-corr-maps-month-diff.png}
  \end{minipage}\hfill
  \begin{minipage}[h]{0.28\linewidth} % Right side for the caption
    \centering \captionof{figure}{Lag correlations as a function of
      start month and lag of Nino-3.4 with temperature tendencies. The
      global mean of the correlation is shown in the title with bold
      indicating the lag with maximum value.
      % White (black) dots indicate statistical significance at the
      % 5\% (1\%) level.
    }\label{fig:corr-with-seasonality-2}
    \begin{enumerate}\scriptsize
      \item 
    \end{enumerate}
  \end{minipage}
\end{frame}

\begin{frame}{Tropical temperatures}
  \begin{minipage}[h]{\linewidth}
    \includegraphics[width=\linewidth]{/Users/tippett/notebooks/NMME/enso-t2m-figures/lag-corr-maps-tropics-month.png}
  \end{minipage}
  
  \begin{minipage}[h]{\linewidth} % Right side for the caption
    \centering \captionof{figure}{Lag correlations as a function of
      start month and lag of Nino-3.4 with temperature. The global
      mean of the correlation is shown in the title with bold
      indicating the lag with maximum value.
      % White (black) dots indicate statistical significance at the
      % 5\% (1\%) level.
    }\label{fig:corr-with-seasonality-2}
    \begin{enumerate}\scriptsize
      \item 
    \end{enumerate}
  \end{minipage}
\end{frame}

\begin{frame}{Tropical temperature tendencies}
  \begin{minipage}[h]{\linewidth}
    \includegraphics[width=\linewidth]{/Users/tippett/notebooks/NMME/enso-t2m-figures/lag-corr-maps-tropics-month-diff.png}
  \end{minipage}
  
  \begin{minipage}[h]{\linewidth} % Right side for the caption
    \centering \captionof{figure}{Lag correlations as a function of
      start month and lag of Nino-3.4 with temperature tendencies. The
      global mean of the correlation is shown in the title with bold
      indicating the lag with maximum value.
      % White (black) dots indicate statistical significance at the
      % 5\% (1\%) level.
    }\label{fig:corr-with-seasonality-2}
    \begin{enumerate}\scriptsize
      \item 
    \end{enumerate}
  \end{minipage}
\end{frame}


\begin{frame}[allowframebreaks]{Tropical and global means: correlation}
  \begin{minipage}[h]{0.65\linewidth}
    \includegraphics[width=\linewidth]{/Users/tippett/notebooks/NMME/enso-t2m-figures/lag-corr-with-seasonality-2.pdf}
  \end{minipage}\hspace{0.5em}
  \begin{minipage}[h]{0.3\linewidth} % Right side for the caption
    \centering \captionof{figure}{Lag correlations as a function of
      start month and lag of Nino-3.4 with (a) tropical mean
      temperature, (b) detrended tropical mean temperature, (c) global
      mean temperature, (d) detrended global mean temperature. White
      (black) dots indicate statistical significance at the 5\% (1\%)
      level. Detrended analysis detrends both time
      series.}\label{fig:corr-with-seasonality-2}
    \begin{enumerate}\scriptsize
      \item 
    \end{enumerate}
  \end{minipage}
\newpage
  \begin{minipage}[h]{0.65\linewidth}
    \includegraphics[width=\linewidth]{/Users/tippett/notebooks/NMME/enso-t2m-figures/lag-corr-with-seasonality-2-target.pdf}
  \end{minipage}\hspace{0.5em}
  \begin{minipage}[h]{0.3\linewidth} % Right side for the caption
    \centering \captionof{figure}{Lag correlations as a function of
      start month and target month of Nino-3.4 with (a) tropical mean
      temperature, (b) detrended tropical mean temperature, (c) global
      mean temperature, (d) detrended global mean temperature. White
      (black) dots indicate statistical significance at the 5\% (1\%)
      level. Detrended analysis detrends both time
      series.}\label{fig:corr-with-seasonality-2}
    \begin{enumerate}\scriptsize
    \item 
    \end{enumerate}
  \end{minipage}
\newpage
  \begin{minipage}[h]{0.65\linewidth}
    \includegraphics[width=\linewidth]{/Users/tippett/notebooks/NMME/enso-t2m-figures/lag-corr-with-seasonality-2-lines.pdf}
  \end{minipage}\hspace{0.5em}
  \begin{minipage}[h]{0.3\linewidth} % Right side for the caption
    \centering \captionof{figure}{Lag correlations as a function of
      start month and target month of Nino-3.4 with (a) tropical mean
      temperature, (b) detrended tropical mean temperature, (c) global
      mean temperature, (d) detrended global mean temperature. Gray
      (black) dots indicate statistical significance at the 5\% (1\%)
      level.}\label{fig:corr-with-seasonality-2}
    \begin{enumerate}\scriptsize
    \item The spread between May and June starts for the tropics is
      interesting. Spring persistence barrier plus something for June
      starts?
    \end{enumerate}
  \end{minipage}
\end{frame}

\begin{frame}[allowframebreaks]{Tropical and global means: regression}
  \begin{minipage}[h]{0.65\linewidth}
    \includegraphics[width=\linewidth]{/Users/tippett/notebooks/NMME/enso-t2m-figures/lag-reg-with-seasonality-2-diff.pdf}
  \end{minipage}\hspace{0.5em}
  \begin{minipage}[h]{0.3\linewidth} % Right side for the caption
    \centering \captionof{figure}{Lag regressions between Nino-3.4
      and tendencies of (a) tropical mean temperature, (b) detrended
      tropical mean temperature, (c) global mean temperature, (d)
      detrended global mean temperature. White (black) dots indicate
      statistical significance at the 5\% (1\%) level. Detrended
      analysis detrends both time
      series.}\label{fig:reg-with-seasonality-2}
    \begin{enumerate}\scriptsize
    \item 
    \end{enumerate}
  \end{minipage}
\newpage
  \begin{minipage}[h]{0.65\linewidth}
    \includegraphics[width=\linewidth]{/Users/tippett/notebooks/NMME/enso-t2m-figures/lag-reg-with-seasonality-2-target.pdf}
  \end{minipage}\hspace{0.5em}
  \begin{minipage}[h]{0.3\linewidth} % Right side for the caption
    \centering \captionof{figure}{Lag regressions between Nino-3.4
      and (a) tropical mean temperature, (b) detrended tropical mean
      temperature, (c) global mean temperature, (d) detrended global
      mean temperature. White (black) dots indicate statistical
      significance at the 5\% (1\%) level. Detrended analysis detrends
      both time series.}\label{fig:reg-with-seasonality-2}
    \begin{enumerate}\scriptsize
    \item 
    \end{enumerate}
  \end{minipage}
\newpage
  \begin{minipage}[h]{0.65\linewidth}
    \includegraphics[width=\linewidth]{/Users/tippett/notebooks/NMME/enso-t2m-figures/lag-reg-with-seasonality-2-lines.pdf}
  \end{minipage}\hspace{0.5em}
  \begin{minipage}[h]{0.3\linewidth} % Right side for the caption
    \centering \captionof{figure}{Lag regression coefficients as a
      function of start month and target month of Nino-3.4 with (a)
      tropical mean temperature, (b) detrended tropical mean
      temperature, (c) global mean temperature, (d) detrended global
      mean temperature. Gray (black) dots indicate statistical
      significance at the 5\% (1\%)
      level.}\label{fig:corr-with-seasonality-2}
    \begin{enumerate}\scriptsize
    \item What's the deal with June?
    \end{enumerate}
  \end{minipage}
\end{frame}

\begin{frame}{Indian and Atlantic sectors}
  \begin{minipage}[h]{0.65\linewidth}
    \includegraphics[width=\linewidth]{/Users/tippett/notebooks/NMME/enso-t2m-figures/lag-corr-with-seasonality-2-target-ind-atl.pdf}
  \end{minipage}\hspace{0.5em}
  \begin{minipage}[h]{0.3\linewidth} % Right side for the caption
    \centering \captionof{figure}{Lag correlations between Nino-3.4
      and (a) Atlantic sector mean temperature, (b) detrended Atlantic
      sector mean temperature, (c) Indian sector mean temperature, (d)
      detrended Indian sector mean temperature. White (black) dots
      indicate statistical significance at the 5\% (1\%)
      level. Detrended analysis detrends both time
      series.}\label{fig:reg-with-seasonality-2}
    \begin{enumerate}\scriptsize
    \item 
    \end{enumerate}
  \end{minipage}
\end{frame}


\begin{frame}{Tropical and global mean tendencies}
  \begin{minipage}[h]{0.65\linewidth}
    \includegraphics[width=\linewidth]{/Users/tippett/notebooks/NMME/enso-t2m-figures/lag-reg-corr-with-seasonality-2-diff-target.pdf}
  \end{minipage}\hspace{0.5em}
  \begin{minipage}[h]{0.3\linewidth} % Right side for the caption
    \centering \captionof{figure}{Lag correlation between Nino-3.4
      and tendencies of (a) tropical mean temperature and (b) global
      mean temperature.  Lag regression between Nino-3.4 and
      tendencies of (c) tropical mean temperature and (d) global mean
      temperature. White (black) dots indicate statistical
      significance at the 5\% (1\%) level.}\label{fig:grainger}
  \end{minipage}
\end{frame}

\begin{frame}
    \begin{enumerate}
    \item Knowing Nino-3.4 tells you that the tropical mean will be
      increasing in fall and decreasing in the following summer
      (Fig. \ref{fig:grainger}).
    \item Knowing Nino-3.4 tells you that the global mean will be
      increasing in Sep, Oct, and Jan.
    \end{enumerate}
\end{frame}

\begin{frame}{One more thing}
  \begin{enumerate}
  \item Something in models?
  \item Prediction something? If you know something now \dots
  \item Previous fitting work with global mean temperature has focused
    on removing ENSO to get a better estimate of the forced trend.
  \item Does knowing the current state of ENSO (or having an ENSO
    forecast) help with predicting future global mean temperature?
  \item Pattern effect?
  \end{enumerate}
\end{frame}

\begin{frame}{Origin story?}
  \begin{minipage}[h]{0.65\linewidth}
    \includegraphics[width=\linewidth]{/Users/tippett/notebooks/NMME/gmt_figures/ACC_n34_gmt_detrended.pdf}
  \end{minipage}\hspace{0.5em}
  \begin{minipage}[h]{0.3\linewidth} % Right side for the caption
    \centering \captionof{figure}{Simultaneous correlation between Nino-3.4 and detrended global mean temperature}
  \end{minipage}
\end{frame}

%\appendix
%\startappendix

\section{Supplemental or ``not shown''}

\renewcommand\thefigure{S\arabic{figure}}
\setcounter{figure}{0}
\renewcommand\thetable{S\arabic{table}}
\setcounter{table}{0}
\renewcommand\theequation{S\arabic{equation}}
\setcounter{equation}{0}

\begin{frame}{A look at the data}
  \begin{minipage}[h]{0.5\linewidth}
    \includegraphics[width=\linewidth]{/Users/tippett/notebooks/NMME/enso-t2m-figures/data-acfs.pdf}
  \end{minipage}\hspace{0.5em}
  \begin{minipage}[h]{0.32\linewidth} % Right side for the caption
    \centering \captionof{figure}{}\label{fig:data-acfs}
  \end{minipage}
\end{frame}


\begin{frame}{Correlation: accounting for seasonality}
  \begin{minipage}[h]{0.4\linewidth}
    \includegraphics[width=\linewidth]{/Users/tippett/notebooks/NMME/enso-t2m-figures/lag-corr-with-seasonality-all.pdf}
  \end{minipage}\hspace{0.5em}
  \begin{minipage}[h]{0.45\linewidth} % Right side for the caption
    \centering \captionof{figure}{Lag correlations between Nino-3.4
      and (a) tropical mean temperature, (b) detrended tropical mean
      temperature, (c) global mean temperature, (d) detrended global
      mean temperature. White (black) dots indicate statistical
      significance at the 5\% (1\%) level. Detrended analysis detrends
      both time series.}\label{fig:corr-with-seasonality-all}
  \end{minipage}
\end{frame}

\begin{frame}{Correlation of differences: accounting for seasonality}
  \begin{minipage}[h]{0.65\linewidth}
    \includegraphics[width=\linewidth]{/Users/tippett/notebooks/NMME/enso-t2m-figures/lag-corr-with-seasonality-2-diff.pdf}
  \end{minipage}\hspace{0.5em}
  \begin{minipage}[h]{0.3\linewidth} % Right side for the caption
    \centering \captionof{figure}{Lag correlations between Nino-3.4
      and differences of (a) tropical mean temperature, (b) detrended tropical mean
      temperature, (c) global mean temperature, (d) detrended global
      mean temperature. White (black) dots indicate statistical
      significance at the 5\% (1\%) level. Detrended analysis detrends
      both time series.}\label{fig:corr-with-seasonality-2-diff}
    \begin{enumerate}\scriptsize
    \item 
    \end{enumerate}
  \end{minipage}
\end{frame}

\begin{frame}{Correlation of differences: accounting for seasonality}
  \begin{minipage}[h]{0.65\linewidth}
    \includegraphics[width=\linewidth]{/Users/tippett/notebooks/NMME/enso-t2m-figures/lag-reg-with-seasonality-2-diff.pdf}
  \end{minipage}\hspace{0.5em}
  \begin{minipage}[h]{0.3\linewidth} % Right side for the caption
    \centering \captionof{figure}{Lag regressions between Nino-3.4
      and differences of (a) tropical mean temperature, (b) detrended tropical mean
      temperature, (c) global mean temperature, (d) detrended global
      mean temperature. White (black) dots indicate statistical
      significance at the 5\% (1\%) level. Detrended analysis detrends
      both time series.}\label{fig:reg-with-seasonality-2-diff}
    \begin{enumerate}\scriptsize
    \item 
    \end{enumerate}
  \end{minipage}
\end{frame}



\begin{frame}[allowframebreaks]{Previous work (tropical Atlantic)}
  \begin{enumerate}
  \item SST not land?
    
  \item Using COADS, \citet{Curtis1995forcing} found that warm
    tropical Pacific conditions (January) are followed by warming a
    few months later (February--May) in the tropical North
    Atlantic. They attributed the delayed warming to weakened North
    Atlantic trades which in turn resulted in reduced latent heat
    flux, downwelling and cloudiness.
    
  \item Using COADS and reconstructed SST, \citet{Enfield1997tropical}
    found highly significant lag (4--5 months) correlations between
    the first EOF of tropical Pacific SST and the second EOF of
    tropical Atlantic SST which describes variability north of the
    ITCZ. The response was found to be strongest in the boreal spring
    and early summer. They argued that the one-two season lag between
    ENSO and Atlantic SST was due to time needed for SST to respond to
    the reduced wind forcing. ``The wind speed changes affect the
    mixed layer temperature primarily through evaporation,
    entrainment, and sensible heat transfer, in decreasing order of
    importance.''  ``This induces a positive anomaly of SST tendency
    that persists for several months, culminating in tropical North
    Atlantic warming in the late spring and early summer, one-two
    seasons after maximum anomalies in the Pacific.''

  \item In an atmospheric modelling study,
    \citet{Saravanan2000interaction} found that ``surface heat flux
    anomalies over the tropical Atlantic associated with ENSO are
    caused by changes in the wind speed as well as by changes in the
    air–sea temperature difference.''

  \item ``The mechanism predicts that the air–sea humidity difference
    variation is a driver of ENSO-related remote tropical surface
    temperature variability, an addition to wind speed and cloudiness
    variations that previous studies have shown to be important''
    \citep{Chiang2002}.
  
  \item \citet{Alexander2002influence} found in NCEP reanalysis and an
    AGCM with a mixed layer model outside the tropical Pacific that
    ``the observed warming in the tropical North Atlantic and cooling
    in the Gulf of Mexico in the winter/spring after ENSO peaks is
    well simulated by the model'' with warming due to latent heat flux
    and shortwave radiation (cloudiness?).
    
  \end{enumerate}
\end{frame}

\begin{frame}[allowframebreaks]{Indian Ocean}
  \begin{enumerate}
  \item SST not land?
  \item ``The tropical IO gradually warms during an ENSO year,
    reaching a maximum during spring (March – May; Figure 13) about
    one season after NINO3 SST has peaked (Figure 12)''
    \citep{Schott2009indian} ``much of this warming is caused by
    ENSO-induced surface changes in surface heat fluxes, particularly
    the wind-induced latent heat and cloud- induced solar radiation
    fluxes.''

  \item ``the lag time is about 3 months for the Indian Ocean, and
    about 5 or 6 months for the tropical North Atlantic and South
    China Sea'' \citep{Klein1999remote}.  ``The remote correlations
    become even more prominent when the ship-observed SSTs and
    continental surface air temperatures lag the ENSO index by 4
    months (Fig.  2c).'' No seasonality.


    ``In an El Ni\~no event, positive SST anomalies usually appear in
    remote ocean basins such as the South China Sea, the Indian Ocean,
    and the tropical North Atlantic approximately 3 to 6 months after
    SST anomalies peak in the tropical Pacific. Ship data from 1952 to
    1992 and satellite data from the 1980s both demonstrate that
    changes in atmospheric circulation accompanying El Ni\~no induce
    changes in cloud cover and evaporation which, in turn, increase
    the net heat flux entering these remote oceans. It is postulated
    that this increased heat flux is responsible for the surface
    warming of these oceans. Specifically, over the eastern Indian
    Ocean and South China Sea, enhanced subsidence during El Ni\~no
    reduces cloud cover and increases the solar radiation absorbed by
    the ocean, thereby leading to enhanced SSTs. In the tropical North
    Atlantic, a weakening of the trade winds during El Ni\~no reduces
    surface evaporation and increases SSTs. These relationships fit
    the concept of an ‘‘atmospheric bridge’’ that connects SST
    anomalies in the central equatorial Pacific to those in remote
    tropical oceans.''

    ``One of the important aspects of ENSO-related teleconnections is
    that they exhibit a strong seasonal dependence.''

    ``In the tropical North Atlantic (5 –22.5 N, 17.5 – 85 W), changes
    in evaporation appear to have a stronger effect than changes in
    cloud cover.'' ``the correlation between cloud cover and the ENSO
    index is negative and rather weak in most months. A notable
    exception is the April–June period''

    ``indicating that the ENSO index is best correlated with both the
    SST tendency and the heat flux at zero time lag.''
  \end{enumerate}
\end{frame}

\end{document}
